% ============================================================================
% MÓDULO 1 - FUNDAMENTOS DE R Y RSTUDIO
% ============================================================================
\chapter{Fundamentos de R y RStudio}
\label{cap:fundamentos}

\begin{objetivos}
\begin{itemize}
  \item Moverte con soltura por RStudio: sus cuatro paneles, los scripts, los
        atajos de teclado y los Proyectos.
  \item Usar R como una calculadora de negocio (IVA, descuentos, márgenes,
        interés compuesto) y guardar resultados en variables.
  \item Distinguir los tipos de datos (números, texto, lógicos, fechas y
        factores) y convertir entre ellos sin caer en las trampas típicas.
  \item Trabajar con vectores, data frames, listas y matrices para responder
        preguntas de negocio.
  \item Automatizar decisiones con \codigo{if}, \codigo{ifelse()} y bucles, y
        crear tus propias funciones de BI.
  \item Instalar y cargar paquetes, importar datos desde CSV y Excel y
        exportar tus resultados.
  \item Resolver un primer análisis de ventas completo, de principio a fin,
        y saber qué hacer cuando R te muestra un error.
\end{itemize}
\end{objetivos}

Imagina que es lunes y acabas de entrar como analista a una cadena de tiendas
minoristas con diez sucursales en México. Tu jefa te comparte un archivo con
las ventas de todo 2023 y te hace cuatro preguntas: \emph{¿cuánto vendimos en
el año?, ¿cuál fue nuestro mejor y nuestro peor mes?, ¿los días con descuento
realmente venden más? y ¿me lo puedes mandar en un Excel?} Con una hoja de
cálculo podrías contestar, pero te tomaría una mañana entera de filtros, tablas
dinámicas y fórmulas copiadas a mano. Y el próximo mes tendrías que repetirlo
todo.

Con R vas a responder esas cuatro preguntas con unas cuantas líneas de código
que puedes guardar, revisar y volver a ejecutar cada mes en segundos. Eso es
lo que hace un analista de \termino{Business Intelligence} (BI, inteligencia
de negocios): convierte datos crudos en respuestas útiles para tomar
decisiones, y lo hace de forma repetible.

Este módulo parte de cero. No necesitas haber programado nunca. Vamos a ir
paso a paso: primero conocerás RStudio, luego aprenderás a guardar y
calcular datos, después a organizarlos en tablas y, al final
(sección~\ref{sec:m1-caso}), resolverás tú mismo el encargo de tu jefa con
datos reales. Cada ejemplo tiene el código comentado, la salida exacta que
imprime R y una explicación de qué significa para el negocio. Te recomiendo
que escribas cada ejemplo en tu computadora en lugar de solo leerlo: programar
se aprende con los dedos.

% ============================================================================
\section{Tu primera sesión en RStudio}
\label{sec:m1-rstudio}
% ============================================================================

\subsection{R y RStudio: el motor y el tablero}

Antes de escribir código conviene distinguir dos programas que siempre
aparecen juntos:

\begin{itemize}
  \item \termino{R} es el \emph{lenguaje} y el motor que hace los cálculos.
        Puede funcionar solo, pero su ventana original es austera.
  \item \termino{RStudio} es un \termino{IDE} (entorno de desarrollo
        integrado): un programa que te da una interfaz cómoda para escribir,
        ejecutar y organizar tu código de R.
\end{itemize}

Piensa en un automóvil: R es el motor y RStudio es el tablero con volante,
velocímetro y espejos. Necesitas los dos instalados, pero en el día a día solo
abrirás RStudio, que a su vez usa R por dentro. Si todavía no los tienes,
sigue las instrucciones de instalación de la introducción del libro
(\emph{Cómo usar este libro}) y vuelve aquí.

\subsection{Los cuatro paneles}

Al abrir RStudio verás una ventana dividida en cuatro zonas llamadas
\termino{paneles}. La figura~\ref{fig:m01-rstudio} muestra su distribución
habitual.

\begin{figure}[H]
\centering
\begin{tikzpicture}[
  font=\sffamily\footnotesize,
  panel/.style={draw=gray!55, fill=white},
  cabecera/.style={fill=gray!14, draw=gray!55},
  pestana/.style={font=\sffamily\scriptsize, anchor=west, inner sep=1pt},
  numero/.style={circle, fill=acento, text=white,
                 font=\sffamily\bfseries\small, inner sep=0pt,
                 minimum size=5.5mm},
  cod/.style={font=\ttfamily\scriptsize, anchor=north west, align=left,
              inner sep=0pt},
  rotulo/.style={font=\sffamily\small\bfseries, text=primario,
                 anchor=south west, align=left},
  nota/.style={font=\sffamily\scriptsize, text=gristexto, anchor=north west,
               align=left, text width=5.6cm},
]
% --- Ventana y barras superiores ------------------------------------------
\draw[rounded corners=4pt, fill=gray!6, draw=gray!70] (0,0) rectangle (15,9.7);
\fill[gray!22, rounded corners=4pt] (0,9.1) rectangle (15,9.7);
\fill[gray!22] (0,9.1) rectangle (15,9.4);
\fill[peligro!75] (0.35,9.4) circle (0.09);
\fill[acento!80] (0.65,9.4) circle (0.09);
\fill[exito!75] (0.95,9.4) circle (0.09);
\node[font=\sffamily\scriptsize] at (7.5,9.4) {curso-bi-rstudio -- RStudio};
\node[anchor=west, font=\sffamily\scriptsize, text=gristexto] at (0.15,8.82)
  {File\quad Edit\quad Code\quad View\quad Plots\quad Session\quad Build%
   \quad Debug\quad Profile\quad Tools\quad Help};
\node[anchor=east, font=\sffamily\scriptsize, text=secundario] at (14.85,8.82)
  {Project: curso-bi-rstudio};
% --- Panel 1: editor de scripts --------------------------------------------
\draw[panel] (0.15,4.55) rectangle (8.45,8.55);
\draw[cabecera] (0.15,8.15) rectangle (8.45,8.55);
\node[pestana] at (0.25,8.35) {\textbf{mi\_primer\_script.R}\ \ \texttt{x}};
\node[draw=gray!55, fill=white, font=\sffamily\scriptsize, inner sep=1.5pt]
  at (7.8,8.35) {$\rightarrow$ Run};
\node[cod, text=gray] at (0.3,7.95) {1\\2\\3};
\node[cod] at (0.65,7.95)
  {{\color{rcomentario}\# Mi primer script}\\
   ventas <- c(120, 95, 140)\\
   mean(ventas)};
\node[rotulo] at (0.3,5.35) {Editor de scripts (Source)};
\node[nota] at (0.3,5.3) {Aquí escribes y guardas tu código.};
\node[numero] at (8.1,4.95) {1};
% --- Panel 2: consola --------------------------------------------------------
\draw[panel] (0.15,0.15) rectangle (8.45,4.35);
\draw[cabecera] (0.15,3.95) rectangle (8.45,4.35);
\node[pestana] at (0.25,4.15) {\textbf{Console}\quad Terminal\quad
  Background Jobs};
\node[cod] at (0.3,3.75)
  {> ventas <- c(120, 95, 140)\\
   > mean(ventas)\\
   {}[1] 118.3333\\
   > \textbar};
\node[rotulo] at (0.3,1.15) {Consola (Console)};
\node[nota] at (0.3,1.1) {R ejecuta las órdenes y muestra los resultados.};
\node[numero] at (8.1,0.55) {2};
% --- Panel 3: entorno ------------------------------------------------------------
\draw[panel] (8.6,4.55) rectangle (14.85,8.55);
\draw[cabecera] (8.6,8.15) rectangle (14.85,8.55);
\node[pestana] at (8.7,8.35) {\textbf{Environment}\quad History\quad
  Connections};
\node[cod, text=gristexto] at (8.75,7.95) {Values};
\node[cod] at (8.75,7.55) {ventas\hspace{3.2em}num [1:3] 120 95 140};
\node[rotulo] at (8.75,5.35) {Entorno (Environment)};
\node[nota] at (8.75,5.3) {Los objetos que has creado y tu historial.};
\node[numero] at (14.5,4.95) {3};
% --- Panel 4: archivos, gráficas, ayuda ----------------------------------------
\draw[panel] (8.6,0.15) rectangle (14.85,4.35);
\draw[cabecera] (8.6,3.95) rectangle (14.85,4.35);
\node[pestana] at (8.7,4.15) {\textbf{Files}\quad Plots\quad Packages\quad
  Help\quad Viewer};
\foreach \y/\t in {3.6/datasets, 3.2/resultados} {
  \fill[acento!45] (8.8,\y-0.12) rectangle (9.05,\y+0.07);
  \node[cod, anchor=west] at (9.2,\y) {\t};
}
\foreach \y/\t in {2.8/curso-bi-rstudio.Rproj, 2.4/mi\_primer\_script.R} {
  \draw[gray!60, fill=white] (8.83,\y-0.13) rectangle (9.02,\y+0.1);
  \node[cod, anchor=west] at (9.2,\y) {\t};
}
\node[rotulo] at (8.75,1.15) {Archivos, gráficas, paquetes y ayuda};
\node[nota] at (8.75,1.1) {Tus carpetas, tus gráficas y la documentación.};
\node[numero] at (14.5,0.55) {4};
\end{tikzpicture}
\caption{Los cuatro paneles de RStudio: (1) editor de scripts, (2) consola,
  (3) entorno e historial y (4) archivos, gráficas, paquetes y ayuda.}
\label{fig:m01-rstudio}
\end{figure}

\begin{enumerate}
  \item \textbf{Editor de scripts} (arriba a la izquierda). Es como un
        documento de Word para tu código: escribes, corriges y guardas. Si
        no lo ves, crea un script nuevo con \menu{File > New File > R Script}.
  \item \textbf{Consola} (abajo a la izquierda). Es donde R \emph{ejecuta}
        las órdenes y responde. El símbolo \codigo{>} se llama
        \termino{prompt} y significa ``estoy listo, dime qué hago''.
  \item \textbf{Entorno} (arriba a la derecha). La pestaña
        \menu{Environment} muestra los datos que has creado (tus
        ``objetos''); \menu{History} guarda las órdenes que has ejecutado.
  \item \textbf{Archivos, gráficas y ayuda} (abajo a la derecha). En
        \menu{Files} navegas por tus carpetas, en \menu{Plots} aparecen las
        gráficas, en \menu{Packages} ves los paquetes instalados y en
        \menu{Help} se abre la documentación.
\end{enumerate}

\begin{consejo}
Puedes cambiar el tema de colores, el tamaño de letra y la distribución de
los paneles en \menu{Tools > Global Options > Appearance} y
\menu{Pane Layout}. Muchos analistas prefieren un tema oscuro para cansar
menos la vista. No afecta en nada a tu código.
\end{consejo}

\subsection{Consola contra script}

Haz clic en la consola, escribe \codigo{2 + 2} y presiona \tecla{Enter}.
R responde de inmediato:

\begin{rcode}
2 + 2
\end{rcode}
\begin{rsalida}
[1] 4
\end{rsalida}

El \codigo{[1]} no es parte del resultado: indica que ese es el
\emph{primer} elemento de la respuesta. Cuando R muestre respuestas largas,
verás al inicio de cada renglón \codigo{[1]}, \codigo{[13]}, etc., para que
sepas en qué posición vas.

La consola es cómoda para pruebas rápidas, pero tiene un problema: lo que
escribes ahí se pierde al cerrar RStudio. Es como hablar por teléfono con R.
El \termino{script}, en cambio, es un archivo de texto con extensión
\archivo{.R} donde escribes tus instrucciones para guardarlas y volver a
ejecutarlas cuando quieras. Es como dejar la receta por escrito.

\begin{concepto}{Regla de oro}
Escribe tu trabajo \emph{siempre} en un script y ejecútalo desde ahí. Usa la
consola solo para consultas rápidas (``¿cuánto vale esta variable?'',
``¿cómo se llama esta columna?''). Un análisis que solo vive en la consola
no se puede revisar, corregir ni repetir el próximo mes.
\end{concepto}

Para ejecutar código desde el script tienes tres opciones:

\begin{itemize}
  \item \textbf{Una línea}: coloca el cursor en cualquier parte de la línea y
        presiona \tecla{Ctrl} + \tecla{Enter} (en Mac, \tecla{Cmd} +
        \tecla{Return}). El cursor baja solo a la siguiente instrucción, así
        que puedes ir presionando el atajo línea por línea.
  \item \textbf{Una selección}: marca varias líneas con el ratón y usa el
        mismo atajo; se ejecuta solo lo seleccionado.
  \item \textbf{Todo el script}: \tecla{Ctrl} + \tecla{Shift} + \tecla{Enter}
        (en Mac, \tecla{Cmd} + \tecla{Shift} + \tecla{Return}), o el botón
        \menu{Source}.
\end{itemize}

El resultado siempre aparece en la consola, precedido de la instrucción que
se ejecutó.

\subsection{Comentarios: notas para tu yo del futuro}

Todo lo que escribas después del símbolo \codigo{\#} es un
\termino{comentario}: R lo ignora por completo. Los comentarios sirven para
explicar \emph{qué} hace el código y, sobre todo, \emph{por qué}. En este
libro todos los ejemplos van comentados en español.

\begin{rcode}
# Esto es un comentario: R no lo ejecuta
120 + 95 + 140   # Ventas de tres días (también puede ir al final)

# ---- Sección: cálculo de IVA -------------------------------------------
1500 * 0.16      # IVA de una venta de $1,500
\end{rcode}
\begin{rsalida}
[1] 355
[1] 240
\end{rsalida}

La línea \codigo{\# ---- Sección: ... ----} es un truco útil: RStudio
reconoce los comentarios que terminan en cuatro guiones o más como
\termino{secciones} y te permite saltar entre ellas con el menú que aparece
en la parte inferior del editor. En scripts largos es una forma muy práctica
de organizarte.

\begin{hazlotu}
Crea un script nuevo (\tecla{Ctrl} + \tecla{Shift} + \tecla{N}), escribe un
comentario con tu nombre y la fecha, y debajo calcula cuánto suman tus tres
últimas compras del súper. Ejecuta la línea con \tecla{Ctrl} +
\tecla{Enter} y observa la consola.
\end{hazlotu}

\subsection{Atajos de teclado que te ahorrarán horas}

No necesitas memorizarlos hoy; vuelve a esta tabla cuando la necesites. Con
el tiempo los usarás sin pensar.

\begin{table}[H]
\centering
\small
\caption{Atajos de teclado más útiles de RStudio.}
\label{tab:m1-atajos}
\begin{tabularx}{\textwidth}{X l l}
\toprule
\textbf{Acción} & \textbf{Windows / Linux} & \textbf{Mac} \\
\midrule
Ejecutar línea o selección & \tecla{Ctrl}+\tecla{Enter} &
  \tecla{Cmd}+\tecla{Return} \\
Ejecutar todo el script & \tecla{Ctrl}+\tecla{Shift}+\tecla{Enter} &
  \tecla{Cmd}+\tecla{Shift}+\tecla{Return} \\
Nuevo script & \tecla{Ctrl}+\tecla{Shift}+\tecla{N} &
  \tecla{Cmd}+\tecla{Shift}+\tecla{N} \\
Guardar el script & \tecla{Ctrl}+\tecla{S} & \tecla{Cmd}+\tecla{S} \\
Escribir la flecha de asignación \codigo{<-} & \tecla{Alt}+\tecla{-} &
  \tecla{Option}+\tecla{-} \\
Escribir el pipe \codigo{\%>\%} & \tecla{Ctrl}+\tecla{Shift}+\tecla{M} &
  \tecla{Cmd}+\tecla{Shift}+\tecla{M} \\
Comentar o descomentar líneas & \tecla{Ctrl}+\tecla{Shift}+\tecla{C} &
  \tecla{Cmd}+\tecla{Shift}+\tecla{C} \\
Autocompletar nombres & \tecla{Tab} & \tecla{Tab} \\
Recuperar la orden anterior (en la consola) & \tecla{$\uparrow$} &
  \tecla{$\uparrow$} \\
Limpiar la consola & \tecla{Ctrl}+\tecla{L} & \tecla{Ctrl}+\tecla{L} \\
Ir al editor / ir a la consola & \tecla{Ctrl}+\tecla{1} / \tecla{Ctrl}+\tecla{2}
  & \tecla{Ctrl}+\tecla{1} / \tecla{Ctrl}+\tecla{2} \\
Cancelar una orden que no termina & \tecla{Esc} & \tecla{Esc} \\
Reiniciar R (sesión limpia) & \tecla{Ctrl}+\tecla{Shift}+\tecla{F10} &
  \tecla{Cmd}+\tecla{Shift}+\tecla{0} \\
Ayuda de la función bajo el cursor & \tecla{F1} & \tecla{F1} \\
Ver todos los atajos & \tecla{Alt}+\tecla{Shift}+\tecla{K} &
  \tecla{Option}+\tecla{Shift}+\tecla{K} \\
\bottomrule
\end{tabularx}
\end{table}

\begin{advertencia}
Si en la consola aparece un signo \codigo{+} en lugar de \codigo{>}, R te
está diciendo ``tu instrucción no está completa, sigo esperando''. Casi
siempre falta cerrar un paréntesis o unas comillas. Presiona \tecla{Esc}
para cancelar, corrige la línea en el script y vuelve a ejecutarla.
\end{advertencia}

\subsection{Guardar tus scripts}

Guarda con \tecla{Ctrl} + \tecla{S}. La primera vez RStudio te pedirá un
nombre: usa nombres descriptivos, sin espacios ni acentos y con la
extensión \archivo{.R}, por ejemplo \archivo{01\_primeros\_pasos.R} o
\archivo{reporte\_ventas\_mensual.R}. Si el nombre de la pestaña del editor
aparece en rojo con un asterisco, hay cambios sin guardar.

\begin{consejo}
Numerar los scripts (\archivo{01\_importar.R}, \archivo{02\_limpiar.R},
\archivo{03\_analizar.R}) deja claro el orden en que deben ejecutarse. Tu
colega (o tú en seis meses) lo agradecerá.
\end{consejo}

\subsection{Proyectos de RStudio y directorio de trabajo}

Cuando le pides a R que abra un archivo, por ejemplo
\archivo{datasets/ventas\_retail.csv}, ¿dónde lo busca? En su
\termino{directorio de trabajo} (\emph{working directory}): la carpeta de tu
computadora que R toma como ``aquí''. Puedes preguntarle cuál es con
\codigo{getwd()} (\emph{get working directory}):

\begin{rcode}
getwd()   # ¿En qué carpeta está trabajando R?
\end{rcode}

R responde con una ruta que depende de tu computadora; en Windows se verá
algo como \codigo{"C:/Users/Ana/Documents/curso-bi-rstudio"} y en Mac algo
como \codigo{"/Users/ana/Documents/curso-bi-rstudio"}. Lo importante es que
esa carpeta sea la del curso, la que contiene \archivo{datasets/}.

La forma profesional de lograrlo son los \termino{Proyectos de RStudio}. Un
proyecto es una carpeta con un archivo \archivo{.Rproj}; al abrirlo, RStudio
fija automáticamente el directorio de trabajo en esa carpeta. Es como tener
un cajón por cliente: todo lo del proyecto (datos, scripts, resultados) vive
junto.

Para convertir la carpeta del curso en un proyecto:

\begin{enumerate}
  \item Ve a \menu{File > New Project > Existing Directory}.
  \item Con \menu{Browse} elige la carpeta \archivo{curso-bi-rstudio} y
        presiona \menu{Create Project}.
  \item RStudio se reinicia dentro del proyecto: verás su nombre arriba a la
        derecha y el archivo \archivo{curso-bi-rstudio.Rproj} en el panel
        \menu{Files}.
\end{enumerate}

A partir de ahora, abre siempre el curso con doble clic en el archivo
\archivo{.Rproj} (o desde \menu{File > Open Project}). La estructura de
carpetas con la que trabajaremos es esta:

\begin{textocodigo}
curso-bi-rstudio/
|-- curso-bi-rstudio.Rproj      <- abre el proyecto con doble clic
|-- datasets/                   <- datos del curso
|   |-- generar_datasets.R
|   |-- ventas_retail.csv
|   |-- clientes.csv
|   |-- ...
|   `-- datos_empresa.xlsx
|-- mis_scripts/                <- aquí guarda tus scripts
`-- resultados/                 <- aquí R guardará lo que exportes
\end{textocodigo}

Si todavía no generas los datos del curso, hazlo ahora \emph{una sola vez}
desde la consola, con el proyecto abierto:

\begin{rnoejecutar}
# Crea los archivos CSV y Excel de la carpeta datasets/ (solo una vez)
source("datasets/generar_datasets.R")
\end{rnoejecutar}

La función \codigo{source()} ejecuta de principio a fin el script que le
indiques. Al terminar verás un mensaje que confirma los siete archivos
creados. Comprueba que R los encuentra:

\begin{rcode}
file.exists("datasets/ventas_retail.csv")   # ¿Existe el archivo?
list.files("datasets")                      # ¿Qué hay en la carpeta?
\end{rcode}
\begin{rsalida}
[1] TRUE
[1] "clientes.csv"       "datos_empresa.xlsx" "empleados.csv"     
[4] "generar_datasets.R" "productos.csv"      "tiendas.csv"       
[7] "transacciones.csv"  "ventas_retail.csv" 
\end{rsalida}

Fíjate en que escribimos \archivo{"datasets/ventas\_retail.csv"} y no la
dirección completa del archivo. Eso es una \termino{ruta relativa}: se
interpreta \emph{a partir} del directorio de trabajo, como cuando dices ``la
segunda puerta a la derecha''. Una \termino{ruta absoluta} es la dirección
completa (\codigo{"C:/Users/Ana/Documents/curso-bi-rstudio/datasets/..."}) y
solo funciona en la computadora de Ana.

\begin{rnoejecutar}
# Cambiar el directorio de trabajo a mano (NO recomendado)
setwd("C:/Users/Ana/Documents/curso-bi-rstudio")        # Windows
setwd("/Users/ana/Documents/curso-bi-rstudio")          # Mac
\end{rnoejecutar}

\begin{advertencia}
Verás \codigo{setwd()} en muchos tutoriales, pero evítalo dentro de tus
scripts: la ruta solo existe en tu computadora, así que el script falla en
cuanto se lo compartes a alguien (o cuando cambias de equipo). Con un
Proyecto de RStudio y rutas relativas, tu código funciona igual en cualquier
computadora.
\end{advertencia}

\begin{hazlotu}
Abre el proyecto del curso, ejecuta \codigo{getwd()} y confirma que la ruta
termina en \archivo{curso-bi-rstudio}. Después crea la carpeta
\archivo{mis\_scripts} desde el panel \menu{Files} (botón \menu{New Folder})
y guarda ahí tu primer script.
\end{hazlotu}

% ============================================================================
\section{R como calculadora de negocio}
\label{sec:m1-calculadora}
% ============================================================================

La forma más sencilla de usar R es como una calculadora muy potente. Los
\termino{operadores aritméticos} son los de siempre, más dos que te serán
útiles con inventarios:

\begin{table}[H]
\centering
\small
\caption{Operadores aritméticos de R.}
\label{tab:m1-operadores}
\begin{tabularx}{\textwidth}{l l X}
\toprule
\textbf{Operador} & \textbf{Significado} & \textbf{Ejemplo de negocio} \\
\midrule
\codigo{+} & Suma & Ventas de enero más ventas de febrero \\
\codigo{-} & Resta & Ingresos menos costos \\
\codigo{*} & Multiplicación & Precio por unidades vendidas \\
\codigo{/} & División & Ventas entre número de tickets \\
\codigo{\textasciicircum} & Potencia & Crecimiento compuesto $(1 + r)^n$ \\
\codigo{\%/\%} & División entera & Cajas completas que salen de un pedido \\
\codigo{\%\%} & Residuo (módulo) & Piezas sueltas que sobran \\
\bottomrule
\end{tabularx}
\end{table}

\begin{rcode}
15000 + 18500      # Ventas de enero más febrero
485000 - 291000    # Ingresos menos costo de lo vendido
349 * 85           # Precio de un mouse por unidades vendidas
103700 / 7         # Venta promedio diaria de una semana
2^10               # Potencia: 2 elevado a la 10
125 %/% 12         # 125 piezas en cajas de 12: cajas completas
125 %% 12          # ... y piezas sueltas que sobran
\end{rcode}
\begin{rsalida}
[1] 33500
[1] 194000
[1] 29665
[1] 14814.29
[1] 1024
[1] 10
[1] 5
\end{rsalida}

Con 125 piezas llenas 10 cajas de 12 y te sobran 5. Ese tipo de cálculo
aparece a diario en almacenes y compras.

\subsection{El orden de las operaciones importa}

R respeta la \termino{precedencia} de operadores que aprendiste en la
escuela: primero paréntesis, luego potencias, después multiplicaciones y
divisiones, y al final sumas y restas. Un paréntesis mal puesto cambia por
completo un resultado de negocio:

\begin{rcode}
# Precio de $1,000 más 16% de IVA
1000 + 1000 * 0.16      # Correcto: primero el IVA, luego la suma
(1000 + 1000) * 0.16    # Incorrecto: suma primero y luego multiplica
1000 * (1 + 0.16)       # Forma compacta y correcta
\end{rcode}
\begin{rsalida}
[1] 1160
[1] 320
[1] 1160
\end{rsalida}

\begin{consejo}
Cuando dudes, usa paréntesis. No cuestan nada y hacen el código más fácil de
leer: \codigo{(ingresos - costos) / ingresos} se entiende al instante.
\end{consejo}

\subsection{Funciones matemáticas}

Una \termino{función} es una orden con nombre que recibe datos entre
paréntesis y devuelve un resultado, igual que \codigo{=REDONDEAR(A1; 2)} en
Excel. Los datos que recibe se llaman \termino{argumentos} y se separan con
comas.

\begin{rcode}
round(37.77778, 1)     # Redondear a 1 decimal
round(16149.15)        # Sin segundo argumento redondea a entero
ceiling(125 / 12)      # Redondear hacia arriba: cajas necesarias
floor(9.99)            # Redondear hacia abajo
sqrt(144)              # Raíz cuadrada
abs(95000 - 102000)    # Valor absoluto: error del pronóstico
log(100)               # Logaritmo natural
log10(1000)            # Logaritmo base 10
exp(1)                 # El número e
\end{rcode}
\begin{rsalida}
[1] 37.8
[1] 16149
[1] 11
[1] 9
[1] 12
[1] 7000
[1] 4.60517
[1] 3
[1] 2.718282
\end{rsalida}

Observa la diferencia entre \codigo{round()}, \codigo{ceiling()} y
\codigo{floor()}. Si necesitas empacar 125 piezas en cajas de 12, no te
sirven 10.42 cajas ni 10 cajas: necesitas 11, así que usas
\codigo{ceiling()}. Y \codigo{abs()} te da el tamaño de un error sin
importar si te pasaste o te quedaste corto: el pronóstico falló por 7,000
pesos.

\subsection{Ejemplos de negocio}

\begin{ejemplo}{IVA y descuentos}
Una tienda vende un producto en \$1,250 antes de impuestos. ¿Cuánto es el
IVA del 16\% y cuánto paga el cliente? Si un ticket ya incluye IVA y suma
\$2,320, ¿cuál era el precio sin IVA? Y si a una laptop de \$18,999 le
aplicas 15\% de descuento, ¿en cuánto queda?
\end{ejemplo}

\begin{rcode}
1250 * 0.16            # IVA de un producto de $1,250
1250 * 1.16            # Total que paga el cliente
2320 / 1.16            # Precio sin IVA de un ticket de $2,320
2320 - 2320 / 1.16     # IVA incluido en ese ticket
18999 * (1 - 0.15)     # Laptop con 15% de descuento
\end{rcode}
\begin{rsalida}
[1] 200
[1] 1450
[1] 2000
[1] 320
[1] 16149.15
\end{rsalida}

El cliente paga \$1,450, de los cuales \$200 son IVA. El ticket de \$2,320
tenía una base de \$2,000 y \$320 de impuesto, y la laptop con descuento
queda en \$16,149.15. Nota que para quitar el IVA \emph{se divide} entre
1.16; restar el 16\% del total sería un error muy común (y caro).

\begin{advertencia}
Los descuentos sucesivos no se suman. Un 20\% seguido de un 10\% adicional
no es un 30\%: \codigo{1000 * 0.80 * 0.90} da 720, es decir, un descuento
efectivo del 28\%. Cuando diseñes promociones, calcula siempre el precio
final.
\end{advertencia}

\begin{ejemplo}{Interés compuesto y crecimiento}
Inviertes \$50,000 al 10\% anual durante 3 años. ¿Cuánto tendrás? ¿Y si la
tasa es del 12\% anual pero se capitaliza cada mes? Por otro lado, las
ventas de tu empresa pasaron de \$1.2 millones en 2020 a \$1.8 millones en
2023: ¿cuál fue su crecimiento anual compuesto?
\end{ejemplo}

\begin{rcode}
50000 * (1 + 0.10)^3           # Capital después de 3 años al 10%
50000 * (1 + 0.12 / 12)^36     # 12% anual capitalizable cada mes
(1800000 / 1200000)^(1 / 3) - 1  # Crecimiento anual compuesto (TCAC)
log(2) / log(1.08)             # Años para duplicar ventas creciendo 8%
\end{rcode}
\begin{rsalida}
[1] 66550
[1] 71538.44
[1] 0.1447142
[1] 9.006468
\end{rsalida}

La inversión llega a \$66,550 con capitalización anual y a \$71,538.44 si se
capitaliza cada mes. Las ventas crecieron en promedio 14.47\% cada año: esa
es la \termino{tasa de crecimiento anual compuesta} (TCAC o
\emph{CAGR} en inglés), un indicador que verás en muchos informes
financieros. Y con un crecimiento del 8\% anual, las ventas se duplican en
unos 9 años.

\begin{ejemplo}{Margen de un producto}
Un producto se vende en \$450 y cuesta \$280. ¿Cuánto ganas por unidad,
cuál es el margen sobre el precio y cuál es el \emph{markup} (ganancia
sobre el costo)? ¿Y cuál es la cantidad económica de pedido si vendes
12,000 unidades al año, cada pedido cuesta \$500 y mantener una unidad en
inventario cuesta \$12 al año?
\end{ejemplo}

\begin{rcode}
450 - 280                        # Ganancia por unidad
(450 - 280) / 450 * 100          # Margen sobre el precio (%)
(450 - 280) / 280 * 100          # Markup sobre el costo (%)
sqrt(2 * 12000 * 500 / 12)       # Cantidad económica de pedido (EOQ)
\end{rcode}
\begin{rsalida}
[1] 170
[1] 37.77778
[1] 60.71429
[1] 1000
\end{rsalida}

Ganas \$170 por unidad: un margen de 37.8\% sobre el precio o un markup de
60.7\% sobre el costo. Son dos formas distintas de expresar la misma
ganancia, y confundirlas es un error clásico en las juntas de precios. La
fórmula de la \termino{cantidad económica de pedido} (EOQ) dice que lo más
barato es pedir 1,000 unidades cada vez.

\begin{hazlotu}
Un cliente compra 3 artículos de \$349, \$1,299 y \$89 (precios sin IVA).
Calcula en R el subtotal, el IVA y el total. Después aplica un cupón del
10\% sobre el subtotal \emph{antes} del IVA y compara. Pista: usa
paréntesis.
\end{hazlotu}

% ============================================================================
\section{Variables: guarda tus resultados con nombre}
\label{sec:m1-variables}
% ============================================================================

Hasta ahora R calcula y olvida. Para reutilizar un valor necesitas guardarlo
en una \termino{variable} (en R también se le llama \termino{objeto}). La
mejor analogía es una celda de Excel a la que le pones nombre: en lugar de
recordar que las ventas de enero están en \codigo{B2}, las llamas
\codigo{ventas\_enero}.

\subsection{Asignación con \texorpdfstring{\codigo{<-}}{<-}}

Para guardar un valor se usa la flecha \codigo{<-}, que se lee ``recibe'':

\begin{rcode}
ventas_enero <- 15000      # ventas_enero "recibe" 15000
ventas_febrero <- 18500
ventas_marzo <- 22000
\end{rcode}

Al ejecutar estas líneas, la consola no muestra nada (asignar no imprime),
pero en el panel \menu{Environment} aparecen las tres variables con su valor.
Para ver el contenido, escribe su nombre o usa \codigo{print()}:

\begin{rcode}
ventas_enero                 # Escribir el nombre muestra el valor
print(ventas_febrero)        # print() hace lo mismo de forma explícita
total_trimestre <- ventas_enero + ventas_febrero + ventas_marzo
total_trimestre
promedio_mensual <- total_trimestre / 3
promedio_mensual
\end{rcode}
\begin{rsalida}
[1] 15000
[1] 18500
[1] 55500
[1] 18500
\end{rsalida}

Ahora las variables funcionan como las celdas de una hoja de cálculo: si
mañana corriges \codigo{ventas\_marzo}, basta volver a ejecutar las líneas
de abajo para que el total y el promedio se actualicen.

\begin{consejo}
El atajo \tecla{Alt} + \tecla{-} (en Mac, \tecla{Option} + \tecla{-})
escribe \codigo{<-} con espacios incluidos. También funciona \codigo{=} para
asignar, pero la comunidad de R usa \codigo{<-} y así lo verás en todos los
manuales. Si quieres asignar e imprimir a la vez, rodea la línea con
paréntesis: \codigo{(meta <- 50000)}.
\end{consejo}

\subsection{Reglas para nombrar variables}

Un buen nombre explica qué hay dentro. R impone algunas reglas y la
comunidad añade algunas convenciones:

\begin{itemize}
  \item Deben empezar con una letra (o un punto) y pueden contener letras,
        números, puntos y guiones bajos: \codigo{ventas\_2023} es válido,
        \codigo{2023\_ventas} no.
  \item No pueden tener espacios ni guiones medios: \codigo{ventas enero} y
        \codigo{ventas-enero} provocan errores (el guion se interpreta como
        una resta).
  \item R distingue mayúsculas de minúsculas: \codigo{Ventas} y
        \codigo{ventas} son dos variables distintas.
  \item No uses palabras reservadas (\codigo{if}, \codigo{for},
        \codigo{TRUE}, \codigo{function}\ldots) ni nombres de funciones
        existentes (\codigo{c}, \codigo{mean}, \codigo{sum}, \codigo{data}).
  \item Evita acentos y la ñ en los nombres de variables
        (\codigo{anio} en lugar de \codigo{año}): funcionan, pero pueden dar
        problemas al abrir el script en otra computadora.
\end{itemize}

En este libro usamos \termino{snake\_case}: palabras en minúsculas unidas
con guion bajo, como \codigo{ventas\_netas\_mensuales}. Es el estilo más
común en R y el que recomienda la guía de estilo del \emph{tidyverse}.

\begin{table}[H]
\centering
\small
\caption{Nombres de variables: buenos y malos.}
\label{tab:m1-nombres}
\begin{tabularx}{\textwidth}{l l X}
\toprule
\textbf{Nombre} & \textbf{¿Sirve?} & \textbf{Comentario} \\
\midrule
\codigo{ventas\_netas\_2023} & \ding{51} Recomendado & Descriptivo y en
  snake\_case \\
\codigo{margen\_pct} & \ding{51} Recomendado & Abreviatura clara \\
\codigo{VentasNetas} & Funciona & Estilo camelCase, poco usado en R \\
\codigo{x}, \codigo{dato2} & Funciona & No dice qué contiene \\
\codigo{ventas netas} & \ding{55} Error & Tiene un espacio \\
\codigo{2023\_ventas} & \ding{55} Error & Empieza con número \\
\codigo{ventas-netas} & \ding{55} Error & R lo lee como ``ventas menos
  netas'' \\
\bottomrule
\end{tabularx}
\end{table}

\subsection{Sobrescribir, listar y borrar variables}

Si asignas un nuevo valor a una variable existente, el valor anterior se
reemplaza sin aviso. Esto es útil para ir acumulando:

\begin{rcode}
caja <- 5000             # Fondo de caja al abrir la tienda
caja <- caja + 1850      # Primera venta del día
caja <- caja + 920       # Segunda venta
caja <- caja - 300       # Pago a un proveedor en efectivo
caja
\end{rcode}
\begin{rsalida}
[1] 7470
\end{rsalida}

La línea \codigo{caja <- caja + 1850} se lee: ``toma el valor actual de
\codigo{caja}, súmale 1,850 y guarda el resultado otra vez en
\codigo{caja}''. Al final del día hay \$7,470 en caja.

Para saber qué variables existen usa \codigo{ls()} y para borrar una,
\codigo{rm()}:

\begin{rcode}
ls()                      # Lista los objetos del entorno
rm(caja)                  # Borra la variable caja
exists("caja")            # ¿Todavía existe?
\end{rcode}
\begin{rsalida}
[1] "caja"             "promedio_mensual" "total_trimestre" 
[4] "ventas_enero"     "ventas_febrero"   "ventas_marzo"    
[1] FALSE
\end{rsalida}

\begin{rnoejecutar}
rm(list = ls())   # Borra TODOS los objetos del entorno (¡cuidado!)
\end{rnoejecutar}

\begin{consejo}
En lugar de borrar todo con \codigo{rm(list = ls())}, es mejor reiniciar R
(\menu{Session > Restart R}) y volver a ejecutar tu script desde el
principio. Así compruebas que tu análisis funciona ``desde cero'' y no
depende de algo que creaste a mano en la consola. Configura también
\menu{Tools > Global Options > General} para que RStudio \emph{no} guarde ni
restaure el archivo \archivo{.RData} al cerrar: tu script es la fuente de
verdad, no la memoria de la sesión.
\end{consejo}

\subsection{Mini caso: los KPIs de una tienda}

Un \termino{KPI} (\emph{Key Performance Indicator}, indicador clave de
desempeño) es un número que resume qué tan bien va el negocio. Vamos a
calcular los KPIs básicos de una sucursal.

\begin{ejemplo}{Tablero de la Sucursal Centro}
En marzo la Sucursal Centro vendió \$485,000. El costo de la mercancía
vendida fue de \$291,000 y los gastos de operación (renta, nómina, luz)
sumaron \$98,500. La meta del mes era de \$520,000. Calcula la utilidad
bruta, la utilidad de operación, los márgenes y el porcentaje de
cumplimiento de la meta.
\end{ejemplo}

\begin{rcode}
# ---- Datos de entrada ----------------------------------------------------
ingresos <- 485000
costo_ventas <- 291000
gastos_operacion <- 98500
meta_ventas <- 520000

# ---- KPIs ----------------------------------------------------------------
utilidad_bruta <- ingresos - costo_ventas
utilidad_operacion <- utilidad_bruta - gastos_operacion
margen_bruto_pct <- utilidad_bruta / ingresos * 100
margen_operacion_pct <- utilidad_operacion / ingresos * 100
cumplimiento_pct <- ingresos / meta_ventas * 100
faltante_meta <- meta_ventas - ingresos

# ---- Resultados ----------------------------------------------------------
utilidad_bruta
utilidad_operacion
round(margen_bruto_pct, 1)
round(margen_operacion_pct, 1)
round(cumplimiento_pct, 1)
faltante_meta
\end{rcode}
\begin{rsalida}
[1] 194000
[1] 95500
[1] 40
[1] 19.7
[1] 93.3
[1] 35000
\end{rsalida}

\textbf{¿Qué nos dice esto?} La sucursal genera \$194,000 de utilidad bruta
(40\% de margen) y, después de pagar sus gastos, le quedan \$95,500 de
utilidad de operación: por cada 100 pesos vendidos gana 19.7. Sin embargo,
solo cumplió el 93.3\% de la meta y le faltaron \$35,000. Si el próximo mes
cambian las cifras, solo actualizas las cuatro líneas de ``Datos de
entrada'' y vuelves a ejecutar todo: ese es el poder de un script.

\begin{hazlotu}
Crea variables con los datos de una tienda imaginaria (ingresos, costo de
ventas, gastos y meta) y calcula los mismos seis KPIs. Después cambia solo
los ingresos a un valor que supere la meta y vuelve a ejecutar: ¿qué pasa
con \codigo{faltante\_meta}? ¿Cómo lo interpretarías?
\end{hazlotu}

% ============================================================================
\section{Tipos de datos}
\label{sec:m1-tipos}
% ============================================================================

En Excel una celda puede contener un número, un texto, una fecha o un
VERDADERO/FALSO, y Excel decide qué operaciones permite. En R pasa lo mismo:
cada valor tiene un \termino{tipo de dato}, y el tipo determina qué puedes
hacer con él. Sumar dos precios tiene sentido; sumar dos nombres de
sucursal, no. La mayoría de los errores de un principiante vienen de un
tipo equivocado, así que vale la pena conocerlos bien.

\begin{table}[H]
\centering
\small
\caption{Tipos de datos básicos de R.}
\label{tab:m1-tipos}
\begin{tabularx}{\textwidth}{l l X}
\toprule
\textbf{Tipo (\codigo{class})} & \textbf{Ejemplo} & \textbf{Uso en negocio} \\
\midrule
\codigo{numeric} & \codigo{349.90}, \codigo{12} & Precios, montos,
  cantidades, porcentajes \\
\codigo{integer} & \codigo{42L} & Conteos (tickets, clientes); la \codigo{L}
  indica entero \\
\codigo{character} & \codigo{"Sucursal Centro"} & Nombres, ciudades,
  códigos de producto \\
\codigo{logical} & \codigo{TRUE}, \codigo{FALSE} & ¿Es premium? ¿Cumplió la
  meta? ¿Tiene descuento? \\
\codigo{Date} & \codigo{as.Date("2023-03-15")} & Fechas de venta, de
  registro, de vencimiento \\
\codigo{factor} & \codigo{factor(c("CH", "M"))} & Categorías con valores
  fijos: tallas, regiones, niveles \\
\bottomrule
\end{tabularx}
\end{table}

\subsection{Números, texto y lógicos}

La función \codigo{class()} te dice de qué tipo es un objeto:

\begin{rcode}
precio <- 349.90                 # numeric: número con decimales
unidades <- 12                   # también numeric, aunque sea entero
tickets <- 42L                   # integer: la L fuerza un entero
sucursal <- "Sucursal Centro"    # character: texto entre comillas
es_premium <- TRUE               # logical: verdadero o falso

class(precio)
class(unidades)
class(tickets)
class(sucursal)
class(es_premium)
\end{rcode}
\begin{rsalida}
[1] "numeric"
[1] "numeric"
[1] "integer"
[1] "character"
[1] "logical"
\end{rsalida}

Tres detalles importantes:

\begin{itemize}
  \item El texto va \emph{siempre} entre comillas (dobles o simples). Sin
        comillas, R cree que \codigo{Centro} es el nombre de una variable y
        la busca.
  \item \codigo{TRUE} y \codigo{FALSE} van en mayúsculas y sin comillas.
        \codigo{"TRUE"} (con comillas) es un texto, no un lógico.
  \item En la práctica casi nunca necesitas escribir la \codigo{L}: R
        trabaja con \codigo{numeric} y lo convierte cuando hace falta.
\end{itemize}

Existe también \codigo{typeof()}, que muestra cómo guarda R el dato por
dentro. La diferencia se entiende con las fechas, como verás en un momento.

\begin{rcode}
typeof(precio)      # "double": número de doble precisión
typeof(tickets)
\end{rcode}
\begin{rsalida}
[1] "double"
[1] "integer"
\end{rsalida}

Los valores lógicos aparecen cada vez que haces una
\termino{comparación}. Los operadores de comparación son \codigo{==}
(igual a; ¡con dos signos!), \codigo{!=} (distinto de), \codigo{>},
\codigo{<}, \codigo{>=} y \codigo{<=}:

\begin{rcode}
ventas_hoy <- 15600
meta_diaria <- 14000
ventas_hoy >= meta_diaria          # ¿Cumplimos la meta de hoy?
sucursal == "Sucursal Norte"       # ¿Es la sucursal Norte?
TRUE + TRUE + FALSE                # TRUE vale 1 y FALSE vale 0
\end{rcode}
\begin{rsalida}
[1] TRUE
[1] FALSE
[1] 2
\end{rsalida}

Que \codigo{TRUE} valga 1 y \codigo{FALSE} valga 0 parece una curiosidad,
pero es un truco muy poderoso: más adelante contarás ``cuántos días se
cumplió la meta'' simplemente sumando comparaciones.

\begin{advertencia}
\codigo{=} y \codigo{==} no son lo mismo. \codigo{x = 5} \emph{asigna} el
valor 5 a \codigo{x}; \codigo{x == 5} \emph{pregunta} si \codigo{x} vale 5 y
responde \codigo{TRUE} o \codigo{FALSE}. Confundirlos es uno de los errores
más frecuentes.
\end{advertencia}

\subsection{Comprobar y convertir tipos}

Las funciones \codigo{is.*()} preguntan si un objeto es de cierto tipo y las
funciones \codigo{as.*()} lo convierten a otro:

\begin{rcode}
is.numeric(precio)             # ¿Es número?
is.character(precio)           # ¿Es texto?
as.character(2023)             # Número -> texto
as.numeric("349.90")           # Texto -> número
as.integer(3.9)                # ¡Cuidado! Trunca, no redondea
as.logical(c("TRUE", "FALSE", "T", "sí"))
\end{rcode}
\begin{rsalida}
[1] TRUE
[1] FALSE
[1] "2023"
[1] 349.9
[1] 3
[1]  TRUE FALSE  TRUE    NA
\end{rsalida}

Fíjate en dos cosas: \codigo{as.integer(3.9)} devuelve 3, porque
\emph{corta} los decimales en lugar de redondear; y \codigo{as.logical()}
no entiende \codigo{"sí"}, así que devuelve \codigo{NA}. \codigo{NA}
significa ``no disponible'' (\emph{Not Available}): es la forma en que R
marca un dato faltante o que no pudo interpretar. Lo estudiarás con calma
en la sección~\ref{sec:m1-vectores}.

\subsection{Tres trampas típicas}

\subsubsection{Trampa 1: números con comas o signo de pesos}

Cuando los datos vienen de un sistema o de un Excel mal exportado, los
montos a veces llegan como texto con separador de miles o con signo de
pesos. R no sabe interpretarlos:

\begin{rnoejecutar}
as.numeric("1,234")            # La coma de miles confunde a R
\end{rnoejecutar}
\begin{rsalida}
[1] NA
Warning message:
NAs introduced by coercion 
\end{rsalida}

R devuelve \codigo{NA} y una \termino{advertencia} (\emph{warning}):
``NAs introducidos por coerción''. Una advertencia no detiene el programa,
pero te avisa que algo salió raro. Si tu R está en español verás el mensaje
traducido; en este libro los mostramos en inglés, que es como los
encontrarás en los foros de ayuda. La solución es quitar primero los
caracteres que estorban con \codigo{gsub()} (que estudiarás en la
sección~\ref{sec:m1-texto}):

\begin{rcode}
as.numeric(gsub(",", "", "1,234"))         # Quita la coma y convierte
monto_texto <- "$12,850.50"
as.numeric(gsub("[$,]", "", monto_texto))  # Quita $ y comas
\end{rcode}
\begin{rsalida}
[1] 1234
[1] 12850.5
\end{rsalida}

Los corchetes en \codigo{"[\$,]"} significan ``cualquiera de estos
caracteres'': R borra todos los signos \codigo{\$} y todas las comas.

\subsubsection{Trampa 2: sumar un número guardado como texto}

Si un número está guardado como texto (entre comillas), R se niega a hacer
aritmética con él:

\begin{rnoejecutar}
precio_texto <- "100"
precio_texto + 1
\end{rnoejecutar}
\begin{rsalida}
Error in precio_texto + 1 : non-numeric argument to binary operator
\end{rsalida}

El mensaje dice ``argumento no numérico en un operador binario'': uno de
los dos lados de la suma no es un número. La solución es convertirlo:
\codigo{as.numeric(precio\_texto) + 1}. Cuando veas este error, revisa con
\codigo{class()} el tipo de tus variables.

\subsubsection{Trampa 3: factores que parecen números}

Un \termino{factor} es el tipo que usa R para \emph{categorías}: variables
que solo pueden tomar unos cuantos valores conocidos, como tallas, regiones
o niveles de cliente. Internamente guarda cada categoría como un número
(1, 2, 3\ldots) y una lista de etiquetas llamadas \termino{niveles}
(\emph{levels}).

\begin{rcode}
tallas <- c("M", "G", "CH", "M", "M", "G", "CH", "EG")
tallas_f <- factor(tallas)            # Niveles en orden alfabético
tallas_f
levels(tallas_f)
table(tallas_f)                       # Conteo por categoría
\end{rcode}
\begin{rsalida}
[1] M  G  CH M  M  G  CH EG
Levels: CH EG G M
[1] "CH" "EG" "G"  "M" 
tallas_f
CH EG  G  M 
 2  1  2  3 
\end{rsalida}

Los niveles quedaron en orden alfabético (CH, EG, G, M), que no es el orden
lógico de las tallas. Para un reporte conviene fijarlo con el argumento
\codigo{levels}:

\begin{rcode}
tallas_f <- factor(tallas, levels = c("CH", "M", "G", "EG"))
table(tallas_f)                       # Ahora en el orden correcto
\end{rcode}
\begin{rsalida}
tallas_f
CH  M  G EG 
 2  3  2  1 
\end{rsalida}

Ahora la trampa. Supón que una columna de precios quedó guardada como factor
(pasaba mucho con versiones antiguas de R al importar un CSV). Si la
conviertes directo a número, obtienes los \emph{códigos internos} de los
niveles, no los precios:

\begin{rcode}
precios_f <- factor(c("150", "80", "150", "200"))
as.numeric(precios_f)                 # ¡MAL! Devuelve los códigos
as.numeric(as.character(precios_f))   # BIEN: primero a texto
\end{rcode}
\begin{rsalida}
[1] 1 3 1 2
[1] 150  80 150 200
\end{rsalida}

El primer resultado (1 3 1 2) no tiene nada que ver con los precios: son las
posiciones de cada valor en la lista de niveles \codigo{"150"},
\codigo{"200"}, \codigo{"80"}. La regla es sencilla: para convertir un
factor a número, pásalo primero a texto con \codigo{as.character()}.

\subsection{Fechas}

Las fechas merecen su propio tipo, \codigo{Date}, porque con ellas se hacen
operaciones especiales: sumar días, calcular antigüedades, extraer el mes o
el día de la semana. La función \codigo{as.Date()} convierte texto en fecha.
Por defecto espera el formato internacional \codigo{"año-mes-día"}:

\begin{rcode}
fecha_pedido <- as.Date("2023-03-15")     # Formato ISO: AAAA-MM-DD
fecha_pedido
class(fecha_pedido)
typeof(fecha_pedido)                      # Por dentro es un número
as.numeric(fecha_pedido)                  # Días desde el 1-ene-1970
\end{rcode}
\begin{rsalida}
[1] "2023-03-15"
[1] "Date"
[1] "double"
[1] 19431
\end{rsalida}

Aquí se ve la diferencia entre \codigo{class()} y \codigo{typeof()}: para ti
y para las funciones de R es una fecha (\codigo{Date}), pero por dentro R la
guarda como el número de días transcurridos desde el 1 de enero de 1970.
Gracias a eso puede restar fechas o sumarles días.

Si tu fecha viene en el formato mexicano \codigo{día/mes/año}, díselo a R
con el argumento \codigo{format}:

\begin{rcode}
fecha_factura <- as.Date("28/02/2023", format = "%d/%m/%Y")
fecha_factura
vencimiento <- fecha_factura + 30          # Crédito a 30 días
vencimiento
fecha_pago <- as.Date("2023-04-10")
fecha_pago - vencimiento                   # ¿Cuántos días de atraso?
as.numeric(fecha_pago - vencimiento)       # Solo el número
\end{rcode}
\begin{rsalida}
[1] "2023-02-28"
[1] "2023-03-30"
Time difference of 11 days
[1] 11
\end{rsalida}

La factura del 28 de febrero vencía el 30 de marzo y se pagó el 10 de
abril: 11 días de atraso. La resta de fechas da un objeto especial llamado
\codigo{difftime} (``diferencia de tiempo''); con \codigo{as.numeric()} te
quedas solo con el número para usarlo en otros cálculos.

La fecha de hoy se obtiene con \codigo{Sys.Date()}. Su resultado, claro,
depende del día en que la ejecutes:

\begin{rcode}
hoy <- Sys.Date()                  # Fecha de hoy según tu computadora
dias_desde_pedido <- as.numeric(hoy - fecha_pedido)
\end{rcode}

Para mostrar una fecha en otro formato se usa \codigo{format()} con los
códigos de la tabla~\ref{tab:m1-formatos-fecha}:

\begin{rcode}
format(fecha_pedido, "%d/%m/%Y")           # Formato mexicano
format(fecha_pedido, "%A %d de %B de %Y")  # Formato largo
format(fecha_pedido, "%Y-%m")              # Año-mes para agrupar
weekdays(fecha_pedido)                     # Día de la semana
months(fecha_pedido)                       # Nombre del mes
\end{rcode}
\begin{rsalida}
[1] "15/03/2023"
[1] "miércoles 15 de marzo de 2023"
[1] "2023-03"
[1] "miércoles"
[1] "marzo"
\end{rsalida}

\begin{table}[H]
\centering
\small
\caption{Códigos de formato de fechas más usados.}
\label{tab:m1-formatos-fecha}
\begin{tabular}{l l l}
\toprule
\textbf{Código} & \textbf{Significado} & \textbf{Ejemplo (15 de marzo de 2023)} \\
\midrule
\codigo{\%d} & Día con dos dígitos & 15 \\
\codigo{\%m} & Mes con dos dígitos & 03 \\
\codigo{\%Y} & Año con cuatro dígitos & 2023 \\
\codigo{\%y} & Año con dos dígitos & 23 \\
\codigo{\%B} & Nombre del mes & marzo \\
\codigo{\%b} & Mes abreviado & mar \\
\codigo{\%A} & Día de la semana & miércoles \\
\bottomrule
\end{tabular}
\end{table}

\begin{advertencia}
Los nombres de días y meses salen en el idioma de tu computadora. Si tu
sistema está en inglés verás \codigo{"Wednesday"} y \codigo{"March"} en
lugar de \codigo{"miércoles"} y \codigo{"marzo"}. Tus cálculos no cambian;
solo las etiquetas.
\end{advertencia}

\begin{ejemplo}{Frecuencia de compra de un cliente}
Un cliente mayorista hizo cinco pedidos en el primer trimestre. ¿Cada
cuántos días compra en promedio? Si mantiene su ritmo, ¿cuándo esperamos
su siguiente pedido? Esta es la base de las alertas de ``cliente en
riesgo'': si pasa mucho más tiempo del habitual sin comprar, alguien de
ventas debería llamarle.
\end{ejemplo}

\begin{rcode}
pedidos <- as.Date(c("2023-01-05", "2023-01-19", "2023-02-09",
                     "2023-02-23", "2023-03-18"))
diff(pedidos)                          # Días entre pedidos consecutivos
dias_promedio <- as.numeric(mean(diff(pedidos)))
dias_promedio
proximo_pedido <- max(pedidos) + dias_promedio
proximo_pedido
format(proximo_pedido, "%A %d/%m/%Y")
\end{rcode}
\begin{rsalida}
Time differences in days
[1] 14 21 14 23
[1] 18
[1] "2023-04-05"
[1] "miércoles 05/04/2023"
\end{rsalida}

El cliente compra cada 18 días en promedio, así que esperamos su siguiente
pedido alrededor del miércoles 5 de abril. La función \codigo{diff()}
calcula la diferencia entre cada elemento y el anterior; la usarás mucho
con ventas diarias o mensuales.

\begin{hazlotu}
Escribe tu fecha de nacimiento con \codigo{as.Date()} y calcula cuántos
días has vivido restándola de \codigo{Sys.Date()}. ¿Qué día de la semana
naciste? Pista: \codigo{format(fecha, "\%A")}.
\end{hazlotu}

% ============================================================================
\section{Vectores: muchos datos a la vez}
\label{sec:m1-vectores}
% ============================================================================

Rara vez analizas un solo número. Normalmente tienes las ventas de cada día,
los precios de cada producto o las edades de cada cliente. Para eso R usa
\termino{vectores}: colecciones ordenadas de valores \emph{del mismo tipo}.
Un vector es como una columna de Excel: una serie de celdas una tras otra.
Y es la pieza con la que se construye todo lo demás en R.

\subsection{Crear vectores}

La función \codigo{c()} (de \emph{combinar}) une valores en un vector:

\begin{rcode}
ventas_semana <- c(12500, 9800, 11200, 13400, 18900, 22300, 15600)
dias <- c("lunes", "martes", "miércoles", "jueves", "viernes",
          "sábado", "domingo")
ventas_semana
dias
length(ventas_semana)      # ¿Cuántos elementos tiene?
\end{rcode}
\begin{rsalida}
[1] 12500  9800 11200 13400 18900 22300 15600
[1] "lunes"     "martes"    "miércoles" "jueves"    "viernes"  
[6] "sábado"    "domingo"  
[1] 7
\end{rsalida}

Como un vector solo admite un tipo, si mezclas números y texto R convierte
\emph{todo} a texto sin avisarte:

\begin{rcode}
c(1500, "2300", 980)       # Un solo texto "contamina" el vector
\end{rcode}
\begin{rsalida}
[1] "1500" "2300" "980" 
\end{rsalida}

Los números ahora aparecen entre comillas: ya no puedes sumarlos. Cuando
una columna numérica de un archivo trae un solo valor raro (por ejemplo,
\codigo{"N/D"}), R suele leer la columna completa como texto. Ahora sabes
por qué.

\subsection{Secuencias y repeticiones}

Escribir valores uno por uno es tedioso. Para patrones regulares existen
atajos:

\begin{rcode}
1:12                                   # Del 1 al 12 (números de mes)
seq(0, 100000, by = 25000)             # De 0 a 100 mil, de 25 en 25 mil
seq(10, 100, length.out = 4)           # 4 valores equiespaciados
rep("Norte", times = 3)                # Repetir un valor
rep(c("Q1", "Q2"), times = 2)          # Repetir el vector completo
rep(c("Q1", "Q2"), each = 2)           # Repetir cada elemento
seq(as.Date("2023-01-01"), by = "month", length.out = 4)  # Fechas
\end{rcode}
\begin{rsalida}
 [1]  1  2  3  4  5  6  7  8  9 10 11 12
[1]      0  25000  50000  75000 100000
[1]  10  40  70 100
[1] "Norte" "Norte" "Norte"
[1] "Q1" "Q2" "Q1" "Q2"
[1] "Q1" "Q1" "Q2" "Q2"
[1] "2023-01-01" "2023-02-01" "2023-03-01" "2023-04-01"
\end{rsalida}

La última línea es muy útil para construir calendarios: el primer día de
cada mes, cada lunes, cada quincena.

\subsection{Operaciones vectorizadas}

Aquí empieza la magia de R. Las operaciones se aplican a \emph{todos} los
elementos del vector a la vez, sin que tengas que arrastrar una fórmula
como en Excel. A esto se le llama \termino{vectorización}.

\begin{ejemplo}{Aplicar un descuento a todo el catálogo}
El área comercial decide aplicar un 15\% de descuento a siete productos y
necesita los precios nuevos, los precios con IVA y el ingreso esperado si
se venden las unidades pronosticadas.
\end{ejemplo}

\begin{rcode}
precios <- c(18999, 349, 1299, 4599, 899, 2499, 159)
precios * 0.85                    # 15% de descuento a los 7 de golpe
round(precios * 0.85 * 1.16, 2)   # Precio con descuento e IVA
unidades <- c(3, 40, 12, 5, 25, 8, 60)
ingresos <- precios * 0.85 * unidades   # Elemento por elemento
ingresos
sum(ingresos)                     # Ingreso total esperado
\end{rcode}
\begin{rsalida}
[1] 16149.15   296.65  1104.15  3909.15   764.15  2124.15   135.15
[1] 18733.01   344.11  1280.81  4534.61   886.41  2464.01   156.77
[1] 48447.45 11866.00 13249.80 19545.75 19103.75 16993.20  8109.00
[1] 137315
\end{rsalida}

Cuando operas dos vectores del mismo largo, R empareja los elementos por
posición: el primer precio con la primera cantidad, el segundo con la
segunda, etc. El ingreso esperado de los siete productos es de \$137,315.

¿Y si los vectores tienen distinto largo? R \emph{recicla} el más corto:
lo repite hasta completar el largo del otro. Esto se llama
\termino{reciclaje}:

\begin{rcode}
c(1000, 2000, 3000, 4000) * c(1, 0.5)   # 1, 0.5, 1, 0.5
\end{rcode}
\begin{rsalida}
[1] 1000 1000 3000 2000
\end{rsalida}

Si el largo del vector mayor no es múltiplo del menor, R hace el cálculo
de todos modos, pero te lanza una advertencia:

\begin{rnoejecutar}
c(1000, 2000, 3000) * c(1, 0.5)         # No es múltiplo: advertencia
\end{rnoejecutar}
\begin{rsalida}
[1] 1000 1000 3000
Warning message:
In c(1000, 2000, 3000) * c(1, 0.5) :
  longer object length is not a multiple of shorter object length
\end{rsalida}

El caso más común de reciclaje es multiplicar por un solo número
(\codigo{precios * 0.85}): el 0.85 se recicla para cada precio. Si R te
advierte que el largo del vector más largo no es múltiplo del más corto,
casi siempre es porque te faltó o te sobró un dato: revísalo.

\subsection{Funciones de resumen}

Para describir un vector con pocos números existen las funciones de
resumen. Primero pondremos nombre a cada elemento con \codigo{names()}, así
los resultados serán más fáciles de leer:

\begin{rcode}
names(ventas_semana) <- dias      # Cada venta con su día
ventas_semana
sum(ventas_semana)                # Total de la semana
mean(ventas_semana)               # Promedio diario
median(ventas_semana)             # Valor central
sd(ventas_semana)                 # Desviación estándar
min(ventas_semana)
max(ventas_semana)
range(ventas_semana)              # Mínimo y máximo juntos
\end{rcode}
\begin{rsalida}
    lunes    martes miércoles    jueves   viernes    sábado   domingo 
    12500      9800     11200     13400     18900     22300     15600 
[1] 103700
[1] 14814.29
[1] 13400
[1] 4451.752
[1] 9800
[1] 22300
[1]  9800 22300
\end{rsalida}

La semana sumó \$103,700, con un promedio de \$14,814.29 por día. La
\termino{mediana} (\$13,400) es el valor que queda justo a la mitad al
ordenar los datos: la mitad de los días vendió menos y la otra mitad más.
Está por debajo del promedio porque el viernes y el sábado ``jalan'' el
promedio hacia arriba. La \termino{desviación estándar} (\$4,451.75) mide qué
tanto se alejan los días del promedio: cuanto más grande, más irregulares
son las ventas.

\begin{rcode}
quantile(ventas_semana)                    # Cuartiles
quantile(ventas_semana, probs = 0.9)       # Percentil 90
cumsum(ventas_semana)                      # Acumulado de la semana
diff(ventas_semana)                        # Cambio respecto al día previo
summary(ventas_semana)                     # Resumen en una línea
\end{rcode}
\begin{rsalida}
   0%   25%   50%   75%  100% 
 9800 11850 13400 17250 22300 
  90% 
20260 
    lunes    martes miércoles    jueves   viernes    sábado   domingo 
    12500     22300     33500     46900     65800     88100    103700 
   martes miércoles    jueves   viernes    sábado   domingo 
    -2700      1400      2200      5500      3400     -6700 
   Min. 1st Qu.  Median    Mean 3rd Qu.    Max. 
   9800   11850   13400   14814   17250   22300 
\end{rsalida}

Los \termino{cuartiles} dividen los datos ordenados en cuatro partes
iguales: el 25\% de los días vendió menos de \$11,850 y el 75\% menos de
\$17,250. El \termino{acumulado} (\codigo{cumsum()}) responde preguntas como
``¿en qué día superamos los \$50,000 de la semana?'' (el jueves ya
llevábamos \$46,900 y el viernes \$65,800). Y \codigo{diff()} muestra que la
mayor subida fue del jueves al viernes (+\$5,500) y la mayor caída del
sábado al domingo (\$6,700 menos). \codigo{summary()} resume casi todo en
una línea.

\subsection{Indexación: sacar elementos de un vector}

Para obtener elementos concretos se usan corchetes \codigo{[ ]}. Es lo que
en BI llamarías ``filtrar''. Hay tres formas de indicar qué quieres.

\subsubsection{Por posición}

\begin{rcode}
ventas_semana[1]              # El primer elemento (R empieza en 1)
ventas_semana[c(5, 6)]        # Posiciones 5 y 6
ventas_semana[1:5]            # De la 1 a la 5: días entre semana
ventas_semana[-7]             # Todos menos el séptimo
\end{rcode}
\begin{rsalida}
lunes 
12500 
viernes  sábado 
  18900   22300 
    lunes    martes miércoles    jueves   viernes 
    12500      9800     11200     13400     18900 
    lunes    martes miércoles    jueves   viernes    sábado 
    12500      9800     11200     13400     18900     22300 
\end{rsalida}

\subsubsection{Por nombre}

\begin{rcode}
ventas_semana["sábado"]
ventas_semana[c("sábado", "domingo")]
\end{rcode}
\begin{rsalida}
sábado 
 22300 
 sábado domingo 
  22300   15600 
\end{rsalida}

\subsubsection{Por condición lógica}

Esta es la forma más usada en análisis de datos. Primero escribes una
comparación, que produce un vector de \codigo{TRUE}/\codigo{FALSE}, y
luego la usas dentro de los corchetes: R se queda solo con los elementos
marcados como \codigo{TRUE}.

\begin{rcode}
meta_diaria <- 14000
ventas_semana > meta_diaria                  # Vector lógico
ventas_semana[ventas_semana > meta_diaria]   # Solo los días que cumplen
\end{rcode}
\begin{rsalida}
    lunes    martes miércoles    jueves   viernes    sábado   domingo 
    FALSE     FALSE     FALSE     FALSE      TRUE      TRUE      TRUE 
viernes  sábado domingo 
  18900   22300   15600 
\end{rsalida}

Para saber \emph{en qué posiciones} se cumple una condición usa
\codigo{which()}, y para localizar el máximo o el mínimo,
\codigo{which.max()} y \codigo{which.min()}:

\begin{rcode}
which(ventas_semana > meta_diaria)   # Posiciones que cumplen
which.max(ventas_semana)             # Posición del mejor día
names(which.max(ventas_semana))      # Solo el nombre del mejor día
names(which.min(ventas_semana))      # Nombre del peor día
\end{rcode}
\begin{rsalida}
viernes  sábado domingo 
      5       6       7 
sábado 
     6 
[1] "sábado"
[1] "martes"
\end{rsalida}

\subsection{Operadores lógicos}

Para combinar condiciones se usan los \termino{operadores lógicos}:

\begin{itemize}
  \item \codigo{\&} (``y''): verdadero solo si \emph{ambas} condiciones lo
        son.
  \item \codigo{|} (``o''): verdadero si \emph{al menos una} lo es.
  \item \codigo{!} (``no''): invierte el resultado.
  \item \codigo{\%in\%} (``está en''): verdadero si el valor pertenece a
        una lista.
\end{itemize}

\begin{rcode}
fin_de_semana <- names(ventas_semana) %in% c("sábado", "domingo")
fin_de_semana
ventas_semana[fin_de_semana]                      # Sábado y domingo
ventas_semana[!fin_de_semana & ventas_semana > 12000]  # Entre semana
ventas_semana[ventas_semana < 10000 | ventas_semana > 20000]  # Extremos
\end{rcode}
\begin{rsalida}
[1] FALSE FALSE FALSE FALSE FALSE  TRUE  TRUE
 sábado domingo 
  22300   15600 
  lunes  jueves viernes 
  12500   13400   18900 
martes sábado 
  9800  22300 
\end{rsalida}

La tercera línea se lee: ``días que \emph{no} son fin de semana \emph{y}
vendieron más de \$12,000''. La cuarta encuentra los días atípicos: los
muy bajos \emph{o} los muy altos.

\subsection{Contar con \texorpdfstring{\codigo{sum()}}{sum()} y \texorpdfstring{\codigo{mean()}}{mean()} de una condición}

Recuerda que \codigo{TRUE} vale 1 y \codigo{FALSE} vale 0. Por lo tanto,
sumar un vector lógico \emph{cuenta} cuántos elementos cumplen la
condición, y promediarlo da la \emph{proporción}:

\begin{rcode}
sum(ventas_semana > meta_diaria)       # ¿Cuántos días cumplieron?
mean(ventas_semana > meta_diaria)      # ¿Qué proporción de días?
\end{rcode}
\begin{rsalida}
[1] 3
[1] 0.4285714
\end{rsalida}

Tres de siete días cumplieron la meta: el 42.9\% de la semana. Este patrón,
\codigo{sum(condición)} y \codigo{mean(condición)}, lo usarás
constantemente: ``¿cuántos clientes son premium?'', ``¿qué porcentaje de
pedidos llegó tarde?''.

\subsection{Datos faltantes: \texorpdfstring{\codigo{NA}}{NA}}

En la vida real faltan datos: el sistema de caja se cayó un día, un cliente
no dio su edad, un producto no tiene costo capturado. R los marca con
\codigo{NA}. Lo importante es que \codigo{NA} es ``contagioso'': cualquier
cálculo que lo incluya da \codigo{NA}, porque R no puede saber el resultado
si le falta un dato.

\begin{rcode}
ventas_tienda <- c(8200, 7650, NA, 9100, 10400, NA, 8800)
mean(ventas_tienda)                   # NA: falta información
mean(ventas_tienda, na.rm = TRUE)     # Ignora los NA
sum(ventas_tienda, na.rm = TRUE)
is.na(ventas_tienda)                  # ¿Dónde hay NA?
sum(is.na(ventas_tienda))             # ¿Cuántos NA hay?
which(is.na(ventas_tienda))           # ¿En qué posiciones?
\end{rcode}
\begin{rsalida}
[1] NA
[1] 8830
[1] 44150
[1] FALSE FALSE  TRUE FALSE FALSE  TRUE FALSE
[1] 2
[1] 3 6
\end{rsalida}

El argumento \codigo{na.rm = TRUE} (\emph{NA remove}) le dice a la función
que ignore los faltantes. El promedio de los cinco días con datos es
\$8,830. Fíjate en que es muy distinto a rellenar los faltantes con cero:
eso bajaría el promedio a \$6,307 y haría creer que la tienda vendió peor.

\begin{advertencia}
No reemplaces \codigo{NA} por 0 sin pensarlo. Un cero dice ``no vendimos
nada''; un \codigo{NA} dice ``no sabemos cuánto vendimos''. Antes de decidir,
averigua \emph{por qué} falta el dato.
\end{advertencia}

\subsection{Ejemplo integrador: la semana de una tienda}

\begin{ejemplo}{Reporte semanal de ventas}
El gerente quiere saber: ¿cuál fue el mejor día?, ¿qué días superaron la
meta de \$14,000?, ¿cómo cambió la venta de un día a otro en porcentaje? y
¿cuáles fueron los tres mejores días?
\end{ejemplo}

\begin{rcode}
# Mejor día y su venta
names(which.max(ventas_semana))
max(ventas_semana)

# Días que superaron la meta
names(ventas_semana)[ventas_semana > meta_diaria]

# Crecimiento porcentual respecto al día anterior
crecimiento_pct <- diff(ventas_semana) / head(ventas_semana, -1) * 100
round(crecimiento_pct, 1)

# Los tres mejores días (ordenados de mayor a menor)
sort(ventas_semana, decreasing = TRUE)[1:3]
\end{rcode}
\begin{rsalida}
[1] "sábado"
[1] 22300
[1] "viernes" "sábado"  "domingo"
   martes miércoles    jueves   viernes    sábado   domingo 
    -21.6      14.3      19.6      41.0      18.0     -30.0 
 sábado viernes domingo 
  22300   18900   15600 
\end{rsalida}

\textbf{¿Qué nos dice esto?} El sábado fue el mejor día (\$22,300) y solo
viernes, sábado y domingo superaron la meta: la tienda depende del fin de
semana. El crecimiento muestra un arranque débil (el martes cayó 21.6\%
respecto al lunes), una subida constante hasta el sábado y una caída del
30\% el domingo. Una decisión razonable sería reforzar promociones de
lunes a jueves en lugar de gastar más en el fin de semana, que ya funciona
solo.

En el código aparecen dos funciones nuevas: \codigo{head(x, -1)} devuelve
todos los elementos menos el último (para dividir cada diferencia entre el
día \emph{anterior}) y \codigo{sort()} ordena los valores; con
\codigo{decreasing = TRUE} los pone de mayor a menor, y \codigo{[1:3]} se
queda con los tres primeros.

\begin{hazlotu}
Crea un vector con las ventas de los 12 meses de un negocio que conozcas (o
inventa cifras) con los nombres \codigo{"ene"}, \codigo{"feb"}\ldots\
\codigo{"dic"}. Calcula el total anual, el mejor mes con
\codigo{names(which.max())}, cuántos meses superaron el promedio y el
crecimiento porcentual de cada mes respecto al anterior.
\end{hazlotu}

% ============================================================================
\section{Trabajar con texto}
\label{sec:m1-texto}
% ============================================================================

En BI el texto está en todas partes: nombres de clientes, ciudades, códigos
de producto, y también los mensajes que tú produces (``La sucursal Norte
vendió \$468,900''). R trae funciones para unir, formatear y limpiar texto.
En programación, a un texto se le llama \termino{cadena} (\emph{string}).

\subsection{Unir textos: \texorpdfstring{\codigo{paste()}}{paste()} y \texorpdfstring{\codigo{paste0()}}{paste0()}}

\begin{rcode}
paste("Sucursal", "Norte")               # Une con un espacio
paste("Q", 1:4, sep = "")                # sep = "" no deja espacio
paste0("SKU-", 1:3)                      # paste0() = paste(sep = "")
paste("ventas", c("enero", "febrero"), "2023", sep = "_")
paste(c("CDMX", "Monterrey", "Puebla"), collapse = ", ")
\end{rcode}
\begin{rsalida}
[1] "Sucursal Norte"
[1] "Q1" "Q2" "Q3" "Q4"
[1] "SKU-1" "SKU-2" "SKU-3"
[1] "ventas_enero_2023"   "ventas_febrero_2023"
[1] "CDMX, Monterrey, Puebla"
\end{rsalida}

\codigo{paste()} también es vectorizada: con \codigo{1:4} produce cuatro
textos a la vez, ideal para generar etiquetas de trimestres, códigos o
nombres de archivo. El argumento \codigo{sep} define qué va \emph{entre} las
piezas y \codigo{collapse} junta todos los elementos de un vector en un
solo texto.

\subsection{Mensajes con formato: \texorpdfstring{\codigo{sprintf()}}{sprintf()}}

\codigo{sprintf()} funciona como una plantilla con huecos. Cada hueco
empieza con \codigo{\%} e indica qué tipo de valor va ahí: \codigo{\%s}
para texto, \codigo{\%d} para enteros y \codigo{\%.2f} para números con 2
decimales. Para escribir el propio signo de porcentaje se usa
\codigo{\%\%}.

\begin{rcode}
ventas_norte <- 468900.5
cumplimiento <- 97.6875
sprintf("La sucursal Norte vendió $%.2f", ventas_norte)
sprintf("Cumplimiento de meta: %.1f%%", cumplimiento)
sprintf("%s: %d tickets", c("Centro", "Norte"), c(1520L, 1204L))
sprintf("SKU-%04d", c(7, 42, 315))       # Rellena con ceros a 4 dígitos
\end{rcode}
\begin{rsalida}
[1] "La sucursal Norte vendió $468900.50"
[1] "Cumplimiento de meta: 97.7%"
[1] "Centro: 1520 tickets" "Norte: 1204 tickets" 
[1] "SKU-0007" "SKU-0042" "SKU-0315"
\end{rsalida}

El último ejemplo es muy práctico para construir códigos de producto o de
cliente con longitud fija (\codigo{SKU-0007}). Nota que
\codigo{sprintf()} no pone separador de miles; para eso está
\codigo{format()}.

\subsection{Números legibles: \texorpdfstring{\codigo{format()}}{format()}}

\begin{rcode}
monto <- 1234567.891
format(monto, big.mark = ",")                  # Separador de miles
format(monto, big.mark = ",", nsmall = 2)      # Con 2 decimales
montos <- c(513916.6, 1407.99, 54056.19)
format(montos, big.mark = ",", nsmall = 2)     # Rellena con espacios
format(montos, big.mark = ",", nsmall = 2, trim = TRUE)
\end{rcode}
\begin{rsalida}
[1] "1,234,568"
[1] "1,234,567.89"
[1] "513,916.60" "  1,407.99" " 54,056.19"
[1] "513,916.60" "1,407.99"   "54,056.19" 
\end{rsalida}

El argumento \codigo{big.mark} pone la coma de miles y \codigo{nsmall}
fija el mínimo de decimales. Con un vector, \codigo{format()} rellena con
espacios para que todos midan lo mismo (útil para alinear columnas en la
consola); con \codigo{trim = TRUE} evitas ese relleno.

\subsection{Formatear montos como moneda}

Un reporte ejecutivo debe mostrar \codigo{\$1,234,567.89}, no
\codigo{1234567.891}. Con R base puedes armarlo combinando
\codigo{round()}, \codigo{format()} y \codigo{paste0()}:

\begin{rcode}
paste0("$", format(round(monto, 2), big.mark = ",", nsmall = 2))
paste0("$", format(round(montos, 2), big.mark = ",", nsmall = 2,
                   trim = TRUE))
\end{rcode}
\begin{rsalida}
[1] "$1,234,567.89"
[1] "$513,916.60" "$1,407.99"   "$54,056.19" 
\end{rsalida}

Funciona, pero es largo. El \paquete{scales} es un \termino{paquete}: un
complemento de R que agrega funciones nuevas (los estudiarás en la
sección~\ref{sec:m1-paquetes}). La notación \codigo{scales::dollar()}
significa ``usa la función \codigo{dollar()} del paquete
\paquete{scales}'':

\begin{rcode}
scales::dollar(monto)                        # Redondea a pesos
scales::dollar(monto, accuracy = 0.01)       # Con centavos
scales::dollar(montos)
scales::dollar(15000, suffix = " MXN")       # Agrega un sufijo
scales::percent(0.9327, accuracy = 0.1)      # Proporción -> porcentaje
scales::comma(1234567)                       # Solo separador de miles
\end{rcode}
\begin{rsalida}
[1] "$1,234,568"
[1] "$1,234,567.89"
[1] "$513,917" "$1,408"   "$54,056" 
[1] "$15,000 MXN"
[1] "93.3%"
[1] "1,234,567"
\end{rsalida}

Observa que \codigo{dollar()} decide solo si muestra centavos: cuando hay
montos grandes, redondea a pesos para que la tabla sea fácil de leer; si
quieres centavos siempre, usa \codigo{accuracy = 0.01}. Y
\codigo{percent()} multiplica por 100 y agrega el signo: recibe una
\emph{proporción} (0.9327), no un porcentaje (93.27). Usarás mucho estas
funciones en las gráficas del Módulo~\ref{cap:visualizacion}.

\begin{consejo}
Si al ejecutar \codigo{scales::dollar()} R responde que no existe el
paquete, todavía no lo tienes instalado. En la
sección~\ref{sec:m1-paquetes} verás cómo instalarlo con
\codigo{install.packages("scales")}.
\end{consejo}

\subsection{Mayúsculas, longitud, extraer y reemplazar}

\begin{rcode}
producto <- "  Laptop HP Pavilion 15  "
toupper(producto)                  # Todo a mayúsculas
tolower("SUCURSAL CENTRO")         # Todo a minúsculas
trimws(producto)                   # Quita espacios al inicio y al final
nchar(trimws(producto))            # Número de caracteres

sku <- "ELEC-2023-0153"            # Categoría-año-consecutivo
substr(sku, 1, 4)                  # Caracteres del 1 al 4
substr(sku, 6, 9)                  # Caracteres del 6 al 9
gsub("-", "", sku)                 # Reemplaza "-" por nada
gsub("Sucursal ", "", c("Sucursal Centro", "Sucursal Norte"))
grepl("Premium", c("Sucursal Premium", "Sucursal Outlet"))
\end{rcode}
\begin{rsalida}
[1] "  LAPTOP HP PAVILION 15  "
[1] "sucursal centro"
[1] "Laptop HP Pavilion 15"
[1] 21
[1] "ELEC"
[1] "2023"
[1] "ELEC20230153"
[1] "Centro" "Norte" 
[1]  TRUE FALSE
\end{rsalida}

\begin{itemize}
  \item \codigo{substr(texto, inicio, fin)} extrae un pedazo por posición:
        perfecto para códigos con estructura fija como los SKU.
  \item \codigo{gsub(buscar, reemplazar, texto)} sustituye \emph{todas}
        las apariciones de un patrón (\codigo{sub()} solo la primera).
  \item \codigo{grepl(patrón, texto)} responde \codigo{TRUE} o
        \codigo{FALSE}: ¿el texto contiene el patrón? Sirve para filtrar,
        por ejemplo, todos los productos que contienen ``Premium''.
\end{itemize}

\begin{ejemplo}{Limpiar nombres de ciudades}
El sistema de registro de clientes permite escribir la ciudad a mano. Al
contar clientes por ciudad aparecen ``CDMX'', `` cdmx'' y ``Cdmx '' como si
fueran ciudades distintas. Es un problema clásico de \termino{calidad de
datos}: el reporte sale mal aunque el cálculo sea correcto.
\end{ejemplo}

\begin{rcode}
ciudades <- c("CDMX", " cdmx", "Cdmx ", "Monterrey", "monterrey ",
              "Puebla")
table(ciudades)                              # Seis "ciudades" distintas
ciudades_limpias <- toupper(trimws(ciudades))
table(ciudades_limpias)                      # Tres ciudades reales
\end{rcode}
\begin{rsalida}
ciudades
      cdmx       CDMX      Cdmx   Monterrey monterrey      Puebla 
         1          1          1          1          1          1 
ciudades_limpias
     CDMX MONTERREY    PUEBLA 
        3         2         1 
\end{rsalida}

Con dos funciones anidadas, \codigo{trimws()} para quitar espacios y
\codigo{toupper()} para unificar mayúsculas, el conteo queda bien: 3
clientes en CDMX, 2 en Monterrey y 1 en Puebla. R ejecuta primero la
función de adentro y le pasa el resultado a la de afuera.

\begin{hazlotu}
Tienes los códigos \codigo{c("ROPA-2023-0012", "HOGA-2022-0450",
"ROPA-2023-0131")}. Extrae la categoría (los 4 primeros caracteres) y el
año de cada uno, y construye un mensaje con \codigo{sprintf()} como
\codigo{"Producto 12 de la categoría ROPA (2023)"}. Pista:
\codigo{as.numeric(substr(x, 11, 14))} te da el consecutivo como número.
\end{hazlotu}

% ============================================================================
\section{Data frames y tibbles: tus hojas de cálculo en R}
\label{sec:m1-dataframes}
% ============================================================================

Un \termino{data frame} es la estructura más importante para BI: una tabla
con filas y columnas, igual que una hoja de Excel. Cada \emph{columna} es un
vector (todos sus valores del mismo tipo) y todas las columnas tienen el
mismo largo. Cada \emph{fila} es una observación: un producto, una venta,
un cliente.

\subsection{Crear un data frame}

\begin{ejemplo}{Catálogo de productos}
Vas a analizar ocho productos de tres categorías. Para cada uno conoces su
nombre, categoría, precio de venta, costo y unidades vendidas en el mes.
Quieres saber cuáles generan más ingreso, cuáles tienen mejor margen y
qué categoría deja más utilidad.
\end{ejemplo}

\begin{rcode}
productos <- data.frame(
  nombre = c("Laptop Pro 14", "Mouse inalámbrico", "Monitor 27",
             "Silla ergonómica", "Escritorio", "Playera básica",
             "Chamarra", "Tenis running"),
  categoria = c("Electrónica", "Electrónica", "Electrónica", "Hogar",
                "Hogar", "Ropa", "Ropa", "Ropa"),
  precio = c(18999, 349, 5499, 3899, 4299, 199, 1299, 1899),
  costo = c(14250, 120, 4100, 2300, 2950, 85, 690, 1050),
  unidades = c(12, 85, 20, 15, 9, 240, 38, 44)
)
productos
\end{rcode}
\begin{rsalida}
             nombre   categoria precio costo unidades
1     Laptop Pro 14 Electrónica  18999 14250       12
2 Mouse inalámbrico Electrónica    349   120       85
3        Monitor 27 Electrónica   5499  4100       20
4  Silla ergonómica       Hogar   3899  2300       15
5        Escritorio       Hogar   4299  2950        9
6    Playera básica        Ropa    199    85      240
7          Chamarra        Ropa   1299   690       38
8     Tenis running        Ropa   1899  1050       44
\end{rsalida}

Cada argumento de \codigo{data.frame()} es una columna: primero el nombre
de la columna, luego \codigo{=} y el vector con sus valores. Los números de
la izquierda (1 a 8) son los nombres de las filas.

\subsection{Explorar un data frame}

Cuando recibes una tabla nueva, lo primero es conocerla: ¿cuántas filas y
columnas tiene?, ¿cómo se llaman las columnas?, ¿de qué tipo es cada una?

\begin{rcode}
head(productos, 3)        # Primeras 3 filas (por defecto, 6)
nrow(productos)           # Número de filas
ncol(productos)           # Número de columnas
dim(productos)            # Filas y columnas
names(productos)          # Nombres de las columnas
\end{rcode}
\begin{rsalida}
             nombre   categoria precio costo unidades
1     Laptop Pro 14 Electrónica  18999 14250       12
2 Mouse inalámbrico Electrónica    349   120       85
3        Monitor 27 Electrónica   5499  4100       20
[1] 8
[1] 5
[1] 8 5
[1] "nombre"    "categoria" "precio"    "costo"     "unidades" 
\end{rsalida}

\begin{rcode}
str(productos)            # Estructura: tipo de cada columna
summary(productos[, c("precio", "costo", "unidades")])
\end{rcode}
\begin{rsalida}
'data.frame':   8 obs. of  5 variables:
 $ nombre   : chr  "Laptop Pro 14" "Mouse inalámbrico" "Monitor 27" "Silla ergonómica" ...
 $ categoria: chr  "Electrónica" "Electrónica" "Electrónica" "Hogar" ...
 $ precio   : num  18999 349 5499 3899 4299 ...
 $ costo    : num  14250 120 4100 2300 2950 ...
 $ unidades : num  12 85 20 15 9 240 38 44
     precio          costo            unidades     
 Min.   :  199   Min.   :   85.0   Min.   :  9.00  
 1st Qu.: 1062   1st Qu.:  547.5   1st Qu.: 14.25  
 Median : 2899   Median : 1675.0   Median : 29.00  
 Mean   : 4555   Mean   : 3193.1   Mean   : 57.88  
 3rd Qu.: 4599   3rd Qu.: 3237.5   3rd Qu.: 54.25  
 Max.   :18999   Max.   :14250.0   Max.   :240.00  
\end{rsalida}

\codigo{str()} (de \emph{structure}) es probablemente la función que más
usarás al empezar un análisis: te dice que hay 8 observaciones
(\emph{obs.}) y 5 variables, el tipo de cada una (\codigo{chr} = texto,
\codigo{num} = número) y sus primeros valores. \codigo{summary()} da el
mínimo, los cuartiles, la media y el máximo de cada columna numérica: el
precio promedio es de \$4,555 pero la mediana es de \$2,899, porque la
laptop sube el promedio.

\begin{rnoejecutar}
View(productos)     # Abre la tabla en una pestaña, como en Excel
\end{rnoejecutar}

\codigo{View()} (con \emph{V} mayúscula) abre el data frame en una pestaña
del editor donde puedes ordenar y filtrar con el ratón. Lo mismo ocurre si
haces clic en el nombre de la tabla dentro del panel \menu{Environment}.
Es útil para mirar, pero recuerda que lo que hagas con el ratón no queda
registrado: los filtros reales se escriben en código.

\subsection{Acceder a columnas, filas y celdas}

El signo \codigo{\$} extrae una columna completa como vector. Los corchetes
aceptan dos posiciones separadas por coma: \codigo{[filas, columnas]}. Si
dejas una posición vacía significa ``todas''.

\begin{rcode}
productos$precio                        # Una columna (vector)
mean(productos$precio)                  # ... y con ella calculas
productos[2, ]                          # Fila 2, todas las columnas
productos[2, "precio"]                  # Una celda: fila 2, col. precio
productos[1:3, c("nombre", "precio")]   # Filas 1 a 3, dos columnas
\end{rcode}
\begin{rsalida}
[1] 18999   349  5499  3899  4299   199  1299  1899
[1] 4555.25
             nombre   categoria precio costo unidades
2 Mouse inalámbrico Electrónica    349   120       85
[1] 349
             nombre precio
1     Laptop Pro 14  18999
2 Mouse inalámbrico    349
3        Monitor 27   5499
\end{rsalida}

\subsection{Filtrar filas}

Filtrar un data frame es igual que filtrar un vector: escribes una
condición y la pones en la posición de las \emph{filas}, antes de la coma.

\begin{rcode}
# Productos con precio mayor a $3,000
productos[productos$precio > 3000, ]

# Ropa con más de 40 unidades vendidas (dos condiciones con &)
productos[productos$categoria == "Ropa" & productos$unidades > 40, ]

# Hogar o Ropa, mostrando solo nombre y precio
productos[productos$categoria %in% c("Hogar", "Ropa"),
          c("nombre", "precio")]
\end{rcode}
\begin{rsalida}
            nombre   categoria precio costo unidades
1    Laptop Pro 14 Electrónica  18999 14250       12
3       Monitor 27 Electrónica   5499  4100       20
4 Silla ergonómica       Hogar   3899  2300       15
5       Escritorio       Hogar   4299  2950        9
          nombre categoria precio costo unidades
6 Playera básica      Ropa    199    85      240
8  Tenis running      Ropa   1899  1050       44
            nombre precio
4 Silla ergonómica   3899
5       Escritorio   4299
6   Playera básica    199
7         Chamarra   1299
8    Tenis running   1899
\end{rsalida}

\begin{advertencia}
No olvides la coma: \codigo{productos[productos\$precio > 3000, ]}. Sin
ella, R interpreta la condición como una selección de \emph{columnas} y
responde con el error \codigo{undefined columns selected}. La coma separa
``qué filas'' de ``qué columnas''.
\end{advertencia}

Una alternativa más legible es \codigo{subset()}, donde puedes escribir
los nombres de las columnas sin repetir \codigo{productos\$}:

\begin{rcode}
subset(productos, precio < 1000, select = c(nombre, precio, unidades))
\end{rcode}
\begin{rsalida}
             nombre precio unidades
2 Mouse inalámbrico    349       85
6    Playera básica    199      240
\end{rsalida}

En el Módulo~\ref{cap:manipulacion} aprenderás \codigo{filter()} del
paquete \paquete{dplyr}, que lleva esta idea todavía más lejos.

\subsection{Ordenar}

La función \codigo{order()} devuelve las \emph{posiciones} de las filas en
el orden deseado; al ponerla en la posición de las filas, reordenas la
tabla:

\begin{rcode}
# De menor a mayor precio
productos[order(productos$precio), c("nombre", "precio")]

# Los 3 productos más caros (de mayor a menor)
caros <- productos[order(productos$precio, decreasing = TRUE), ]
caros[1:3, c("nombre", "categoria", "precio")]
\end{rcode}
\begin{rsalida}
             nombre precio
6    Playera básica    199
2 Mouse inalámbrico    349
7          Chamarra   1299
8     Tenis running   1899
4  Silla ergonómica   3899
5        Escritorio   4299
3        Monitor 27   5499
1     Laptop Pro 14  18999
         nombre   categoria precio
1 Laptop Pro 14 Electrónica  18999
3    Monitor 27 Electrónica   5499
5    Escritorio       Hogar   4299
\end{rsalida}

\subsection{Crear columnas calculadas}

Para agregar una columna basta asignarle un vector con \codigo{\$}. Como las
operaciones son vectorizadas, el cálculo se hace para todas las filas a la
vez:

\begin{rcode}
productos$ingreso <- productos$precio * productos$unidades
productos$utilidad <- (productos$precio - productos$costo) *
  productos$unidades
productos$margen_pct <- round((productos$precio - productos$costo) /
                                productos$precio * 100, 1)
productos[, c("nombre", "ingreso", "utilidad", "margen_pct")]
\end{rcode}
\begin{rsalida}
             nombre ingreso utilidad margen_pct
1     Laptop Pro 14  227988    56988       25.0
2 Mouse inalámbrico   29665    19465       65.6
3        Monitor 27  109980    27980       25.4
4  Silla ergonómica   58485    23985       41.0
5        Escritorio   38691    12141       31.4
6    Playera básica   47760    27360       57.3
7          Chamarra   49362    23142       46.9
8     Tenis running   83556    37356       44.7
\end{rsalida}

\textbf{¿Qué nos dice esto?} La laptop es el producto que más ingreso
(\$227,988) y más utilidad (\$56,988) genera, pero con el margen más bajo
(25\%). En cambio, el mouse y la playera, que cuestan poco, dejan márgenes
de 65.6\% y 57.3\%. Un producto puede ser muy importante en ventas y poco
rentable por unidad: por eso en BI siempre se miran juntas varias
métricas.

\subsection{Resumir por grupos}

La pregunta más frecuente en BI es ``¿cuánto por cada\ldots?'': ventas por
tienda, clientes por ciudad, utilidad por categoría. En Excel harías una
tabla dinámica. En R base tienes tres herramientas:

\begin{rcode}
table(productos$categoria)                       # Conteo por grupo
tapply(productos$ingreso, productos$categoria, sum)   # Suma por grupo
aggregate(ingreso ~ categoria, data = productos, FUN = sum)
\end{rcode}
\begin{rsalida}

Electrónica       Hogar        Ropa 
          3           2           3 
Electrónica       Hogar        Ropa 
     367633       97176      180678 
    categoria ingreso
1 Electrónica  367633
2       Hogar   97176
3        Ropa  180678
\end{rsalida}

\begin{itemize}
  \item \codigo{table()} cuenta cuántas filas hay de cada categoría.
  \item \codigo{tapply(valores, grupos, función)} aplica la función a los
        valores de cada grupo y devuelve un vector con nombres.
  \item \codigo{aggregate()} hace lo mismo pero devuelve un data frame,
        ideal para seguir trabajando o exportarlo. La expresión
        \codigo{ingreso \textasciitilde{} categoria} es una
        \termino{fórmula} y se lee ``ingreso según categoría''.
\end{itemize}

\codigo{aggregate()} puede resumir varias columnas a la vez con
\codigo{cbind()}. Con eso armamos una pequeña tabla dinámica de
rentabilidad:

\begin{rcode}
resumen_categoria <- aggregate(cbind(ingreso, utilidad) ~ categoria,
                               data = productos, FUN = sum)
resumen_categoria$margen_pct <- round(resumen_categoria$utilidad /
                                        resumen_categoria$ingreso * 100, 1)
resumen_categoria[order(resumen_categoria$margen_pct,
                        decreasing = TRUE), ]
\end{rcode}
\begin{rsalida}
    categoria ingreso utilidad margen_pct
3        Ropa  180678    87858       48.6
2       Hogar   97176    36126       37.2
1 Electrónica  367633   104433       28.4
\end{rsalida}

\textbf{¿Qué nos dice esto?} Electrónica vende más del doble que Ropa
(\$367,633 contra \$180,678), pero Ropa convierte casi la mitad de cada peso
vendido en utilidad (48.6\%) mientras Electrónica solo el 28.4\%. Si el
equipo comercial tiene presupuesto para una campaña, impulsar Ropa deja
más utilidad por cada peso vendido.

\subsection{Data frame o tibble}

El \termino{tibble} es una versión moderna del data frame que usa el
\emph{tidyverse}, el conjunto de paquetes que estudiarás a partir del
Módulo~\ref{cap:manipulacion}. Se crea con \codigo{tibble::tibble()} y se
comporta casi igual, con algunas mejoras:

\begin{rcode}
productos_tb <- tibble::tibble(
  nombre = c("Laptop Pro 14", "Mouse inalámbrico", "Monitor 27"),
  precio = c(18999, 349, 5499),
  unidades = c(12, 85, 20),
  ingreso = precio * unidades     # Usa columnas recién creadas
)
productos_tb
tibble::as_tibble(productos)      # Convertir un data frame a tibble
\end{rcode}
\begin{rsalida}
# A tibble: 3 x 4
  nombre            precio unidades ingreso
  <chr>              <dbl>    <dbl>   <dbl>
1 Laptop Pro 14      18999       12  227988
2 Mouse inalámbrico    349       85   29665
3 Monitor 27          5499       20  109980
# A tibble: 8 x 8
  nombre   categoria precio costo unidades ingreso utilidad margen_pct
  <chr>    <chr>      <dbl> <dbl>    <dbl>   <dbl>    <dbl>      <dbl>
1 Laptop ~ Electrón~  18999 14250       12  227988    56988       25  
2 Mouse i~ Electrón~    349   120       85   29665    19465       65.6
3 Monitor~ Electrón~   5499  4100       20  109980    27980       25.4
4 Silla e~ Hogar       3899  2300       15   58485    23985       41  
5 Escrito~ Hogar       4299  2950        9   38691    12141       31.4
6 Playera~ Ropa         199    85      240   47760    27360       57.3
7 Chamarra Ropa        1299   690       38   49362    23142       46.9
8 Tenis r~ Ropa        1899  1050       44   83556    37356       44.7
\end{rsalida}

Las diferencias más visibles son:

\begin{itemize}
  \item Al imprimirlo te dice su tamaño (\codigo{8 x 8}) y el tipo de cada
        columna debajo del nombre: \codigo{<chr>} texto, \codigo{<dbl>}
        número decimal, \codigo{<int>} entero, \codigo{<lgl>} lógico,
        \codigo{<date>} fecha.
  \item Solo muestra las primeras 10 filas y las columnas que caben en la
        pantalla; las demás las lista al final. Con tablas de un millón de
        filas, eso evita que la consola se inunde.
  \item Muestra los números con 3 cifras significativas para ahorrar
        espacio (los decimales se atenúan en pantalla), pero el valor
        guardado no cambia.
  \item Dentro de \codigo{tibble()} una columna puede usar otra recién
        creada, como hicimos con \codigo{ingreso}.
  \item Nunca convierte texto en factor y no tiene nombres de fila.
\end{itemize}

Todo lo que aprendiste en esta sección (\codigo{\$}, corchetes,
\codigo{str()}, \codigo{aggregate()}\ldots) funciona igual con tibbles.
Cuando leas archivos con \paquete{readr} o \paquete{readxl} obtendrás
tibbles.

\begin{hazlotu}
Agrega al data frame \codigo{productos} una columna
\codigo{precio\_con\_iva} (precio por 1.16, redondeado a 2 decimales).
Después muestra solo los productos de Electrónica ordenados de mayor a
menor margen y calcula con \codigo{aggregate()} las unidades totales por
categoría.
\end{hazlotu}

% ============================================================================
\section{Listas y matrices}
\label{sec:m1-listas-matrices}
% ============================================================================

Los vectores y los data frames cubren el 90\% de lo que harás en BI, pero
hay dos estructuras más que conviene conocer porque aparecen a menudo. La
figura~\ref{fig:m01-estructuras} resume las cuatro.

\begin{figure}[H]
\centering
\begin{tikzpicture}[
  font=\sffamily\footnotesize,
  celda/.style={draw=gray!60, minimum width=9mm, minimum height=5mm,
                inner sep=0pt, font=\ttfamily\scriptsize},
  titulo/.style={font=\sffamily\small\bfseries, text=primario},
  expl/.style={font=\sffamily\scriptsize, text=gristexto, align=center,
               text width=6.6cm},
]
% ---- Vector --------------------------------------------------------------
\node[titulo] at (3.4,6.6) {Vector};
\foreach \v [count=\i] in {12500, 9800, 11200, 13400, 18900} {
  \node[celda, fill=secundario!12] at (0.9+\i*0.95-0.45,5.95) {\v};
}
\node[expl] at (3.4,5.2) {Una fila de valores, todos del mismo tipo.\\
  \texttt{c(12500, 9800, ...)}};
% ---- Matriz --------------------------------------------------------------
\node[titulo] at (11.2,6.6) {Matriz};
\foreach \fila [count=\r] in {{150,200,180,220},{430,390,510,480},
                              {210,240,190,260}} {
  \foreach \v [count=\c] in \fila {
    \node[celda, fill=morado!10] at (9.3+\c*0.95-0.45,6.45-\r*0.5) {\v};
  }
}
\node[expl] at (11.2,4.4) {Tabla de un solo tipo (normalmente números).\\
  \texttt{matrix(..., nrow = 3)}};
% ---- Data frame -----------------------------------------------------------
\node[titulo] at (3.4,3.4) {Data frame / tibble};
\node[celda, fill=gray!20, minimum width=2.4cm] at (1.5,2.75)
  {\textbf{nombre} <chr>};
\node[celda, fill=gray!20, minimum width=1.7cm] at (3.55,2.75)
  {\textbf{precio} <dbl>};
\node[celda, fill=gray!20, minimum width=1.7cm] at (5.25,2.75)
  {\textbf{premium} <lgl>};
\foreach \n/\p/\l [count=\r] in {Laptop/18999/TRUE, Mouse/349/FALSE,
                                  Silla/3899/TRUE} {
  \node[celda, fill=acento!10, minimum width=2.4cm] at (1.5,2.75-\r*0.5) {\n};
  \node[celda, fill=secundario!12, minimum width=1.7cm] at (3.55,2.75-\r*0.5)
    {\p};
  \node[celda, fill=exito!12, minimum width=1.7cm] at (5.25,2.75-\r*0.5) {\l};
}
\node[expl] at (3.4,0.55) {Columnas de distinto tipo, todas del mismo
  largo.\\ Es la ``hoja de Excel'' de R.};
% ---- Lista ---------------------------------------------------------------
\node[titulo] at (11.2,3.4) {Lista};
\draw[rounded corners=3pt, draw=gray!60, fill=gray!4] (8.4,1.05) rectangle
  (14.0,3.05);
\node[celda, fill=acento!10, minimum width=2.6cm, anchor=west] at (8.6,2.7)
  {nombre: "María"};
\node[celda, fill=exito!12, minimum width=2.2cm, anchor=west] at (11.55,2.7)
  {premium: TRUE};
\node[anchor=west, font=\ttfamily\scriptsize] at (8.6,2.1) {compras:};
\foreach \v [count=\i] in {1250, 830, 2410} {
  \node[celda, fill=secundario!12, minimum width=8mm] at (9.9+\i*0.85,2.1)
    {\v};
}
\node[anchor=west, font=\ttfamily\scriptsize] at (8.6,1.5) {historial:};
\foreach \c in {0,1,2} {
  \foreach \r in {0,1} {
    \node[draw=gray!60, fill=morado!10, minimum width=4mm,
          minimum height=2.5mm, inner sep=0pt] at (10.75+\c*0.42,1.62-\r*0.27) {};
  }
}
\node[anchor=west, font=\sffamily\scriptsize, text=gristexto] at (12.1,1.5)
  {(una tabla)};
\node[expl] at (11.2,0.55) {Una ``carpeta'' que guarda objetos de cualquier
  tipo\\ y tamaño, cada uno con su nombre.};
\end{tikzpicture}
\caption{Las cuatro estructuras de datos básicas de R.}
\label{fig:m01-estructuras}
\end{figure}

\subsection{Matrices: tablas de un solo tipo}

Una \termino{matriz} es una tabla rectangular donde \emph{todos} los
valores son del mismo tipo, casi siempre números. Es como una tabla
dinámica ya resumida: productos en las filas, meses en las columnas.

\begin{ejemplo}{Ventas en unidades por producto y mes}
Tienes las unidades vendidas de tres productos en los primeros cuatro meses
del año. Quieres el total por producto, el total por mes y un pronóstico
para el siguiente cuatrimestre con un crecimiento del 10\%.
\end{ejemplo}

\begin{rcode}
ventas_mat <- matrix(
  c(150, 200, 180, 220,     # Laptop
    430, 390, 510, 480,     # Mouse
    210, 240, 190, 260),    # Monitor
  nrow = 3, byrow = TRUE,   # 3 filas, llenando renglón por renglón
  dimnames = list(c("Laptop", "Mouse", "Monitor"),     # Filas
                  c("ene", "feb", "mar", "abr"))       # Columnas
)
ventas_mat
ventas_mat["Mouse", "mar"]      # Una celda: [fila, columna]
ventas_mat[, "abr"]             # Una columna completa
rowSums(ventas_mat)             # Total por producto
colSums(ventas_mat)             # Total por mes
round(rowMeans(ventas_mat), 1)  # Promedio mensual por producto
ventas_mat * 1.10               # Pronóstico: +10% en todas las celdas
\end{rcode}
\begin{rsalida}
        ene feb mar abr
Laptop  150 200 180 220
Mouse   430 390 510 480
Monitor 210 240 190 260
[1] 510
 Laptop   Mouse Monitor 
    220     480     260 
 Laptop   Mouse Monitor 
    750    1810     900 
ene feb mar abr 
790 830 880 960 
 Laptop   Mouse Monitor 
  187.5   452.5   225.0 
        ene feb mar abr
Laptop  165 220 198 242
Mouse   473 429 561 528
Monitor 231 264 209 286
\end{rsalida}

El mouse es el producto de mayor volumen (1,810 unidades en cuatro meses)
y abril fue el mes más fuerte (960 unidades). Las funciones
\codigo{rowSums()} y \codigo{colSums()} son los ``totales de fila'' y
``totales de columna'' de una tabla dinámica. Y como las matrices también
son vectorizadas, \codigo{ventas\_mat * 1.10} aplica el crecimiento a las
doce celdas de golpe.

\subsection{Listas: una carpeta con objetos de cualquier tipo}

Una \termino{lista} es la estructura más flexible: guarda objetos de
distinto tipo y tamaño, cada uno con su nombre. Piensa en el expediente de
un cliente: una carpeta con su nombre (texto), su fecha de alta (fecha), si
es premium (lógico) y el historial de sus compras (un vector de montos).

\begin{rcode}
cliente <- list(
  id = 1045,
  nombre = "María González",
  ciudad = "Guadalajara",
  premium = TRUE,
  fecha_alta = as.Date("2021-06-14"),
  compras = c(1250.00, 830.50, 2410.00, 675.25)
)
str(cliente)                   # Estructura de la lista
cliente$nombre                 # Un elemento por nombre
cliente[["compras"]]           # Otra forma: doble corchete
sum(cliente$compras)           # Total comprado
mean(cliente$compras)          # Ticket promedio
cliente$canal <- "Tienda en línea"   # Agregar un elemento nuevo
names(cliente)
\end{rcode}
\begin{rsalida}
List of 6
 $ id        : num 1045
 $ nombre    : chr "María González"
 $ ciudad    : chr "Guadalajara"
 $ premium   : logi TRUE
 $ fecha_alta: Date[1:1], format: "2021-06-14"
 $ compras   : num [1:4] 1250 830 2410 675
[1] "María González"
[1] 1250.00  830.50 2410.00  675.25
[1] 5165.75
[1] 1291.438
[1] "id"         "nombre"     "ciudad"     "premium"    "fecha_alta"
[6] "compras"    "canal"     
\end{rsalida}

María ha comprado \$5,165.75 en cuatro visitas, con un ticket promedio de
\$1,291.44. Con \codigo{\$} o con \codigo{[[ ]]} sacas un elemento tal cual
(aquí, el vector de compras) y puedes operar con él.

\begin{concepto}{¿Por qué importan las listas?}
Un data frame es, en realidad, una lista de vectores del mismo largo: por
eso \codigo{productos\$precio} funciona igual que \codigo{cliente\$nombre}.
Además, muchas funciones de R devuelven listas: los resultados de un modelo
estadístico, la respuesta de una API web o la configuración de un reporte.
En el Módulo~\ref{cap:ml} extraerás información de listas así.
\end{concepto}

\begin{hazlotu}
Crea una lista \codigo{tienda} con el nombre de una sucursal, su ciudad,
su fecha de apertura, su número de empleados y un vector con sus ventas de
los últimos seis meses. Calcula la venta promedio y cuántos años lleva
abierta la tienda. Después crea una matriz de 2 sucursales $\times$ 3 meses
y obtén el total por sucursal con \codigo{rowSums()}.
\end{hazlotu}

% ============================================================================
\section{Control de flujo: decisiones y repeticiones}
\label{sec:m1-control}
% ============================================================================

Hasta ahora R ejecuta las instrucciones una tras otra, siempre igual. El
\termino{control de flujo} te permite decirle ``haz esto \emph{solo si}
pasa tal cosa'' o ``repite esto para cada sucursal''. Son las mismas reglas
de negocio que aplicas mentalmente todos los días.

\subsection{Decidir con \texorpdfstring{\codigo{if}}{if} y \texorpdfstring{\codigo{else}}{else}}

La estructura \codigo{if (condición) \{ \ldots \} else \{ \ldots \}}
ejecuta el primer bloque si la condición es \codigo{TRUE} y el segundo si
es \codigo{FALSE}. Las llaves \codigo{\{ \}} agrupan las instrucciones de
cada caso.

\begin{ejemplo}{Clasificar una venta}
La política comercial dice: una venta de \$5,000 o más es ``Grande'' y la
revisa un gerente; de \$1,000 a \$4,999 es ``Mediana''; y por debajo de
\$1,000 es ``Chica''. Queremos que R clasifique una venta de \$3,250.
\end{ejemplo}

\begin{rcode}
monto_venta <- 3250

if (monto_venta >= 5000) {
  tipo_venta <- "Grande"
} else if (monto_venta >= 1000) {
  tipo_venta <- "Mediana"
} else {
  tipo_venta <- "Chica"
}
tipo_venta
\end{rcode}
\begin{rsalida}
[1] "Mediana"
\end{rsalida}

R revisa las condiciones de arriba hacia abajo y se detiene en la primera
que se cumple. Como \$3,250 no llega a 5,000 pero sí supera 1,000, la venta
queda como ``Mediana''. Nota la sangría (los dos espacios al inicio de las
líneas internas): R no la necesita, pero hace el código mucho más fácil de
leer. RStudio la pone sola; si se desordena, selecciona el código y
presiona \tecla{Ctrl} + \tecla{I}.

\begin{advertencia}
\codigo{if} evalúa \emph{un solo} valor lógico. Si le das un vector
(por ejemplo, las ventas de toda la semana), R se detiene con el error
\codigo{the condition has length > 1}. Para evaluar muchos valores a la
vez usa \codigo{ifelse()}, que verás enseguida.
\end{advertencia}

\subsection{Decidir para todo un vector: \texorpdfstring{\codigo{ifelse()}}{ifelse()}}

\codigo{ifelse(condición, valor\_si\_TRUE, valor\_si\_FALSE)} es la versión
vectorizada de \codigo{if}: evalúa la condición elemento por elemento y
devuelve un vector del mismo largo. Es el equivalente directo de la
función \codigo{=SI()} de Excel arrastrada a toda una columna.

\begin{rcode}
ventas_semana
etiqueta_meta <- ifelse(ventas_semana >= meta_diaria,
                        "Cumplió", "No cumplió")
etiqueta_meta

# Anidar ifelse() para tres categorías
nivel_dia <- ifelse(ventas_semana >= 18000, "Alto",
                    ifelse(ventas_semana >= 12000, "Medio", "Bajo"))
table(nivel_dia)
\end{rcode}
\begin{rsalida}
    lunes    martes miércoles    jueves   viernes    sábado   domingo 
    12500      9800     11200     13400     18900     22300     15600 
       lunes       martes    miércoles       jueves      viernes 
"No cumplió" "No cumplió" "No cumplió" "No cumplió"    "Cumplió" 
      sábado      domingo 
   "Cumplió"    "Cumplió" 
nivel_dia
 Alto  Bajo Medio 
    2     2     3 
\end{rsalida}

Con una sola línea etiquetaste los siete días. Anidar \codigo{ifelse()}
funciona para tres o cuatro categorías; con más se vuelve ilegible y es
mejor usar \codigo{cut()} (que aparece en los ejercicios) o
\codigo{case\_when()}, que estudiarás en el Módulo~\ref{cap:manipulacion}.

\subsection{Repetir con \texorpdfstring{\codigo{for}}{for}}

Un \termino{bucle} \codigo{for} repite un bloque de código una vez por
cada elemento de un vector. La variable del bucle (aquí \codigo{i}) va
tomando cada valor en turno.

\begin{ejemplo}{Reporte de cumplimiento por sucursal}
Tienes las ventas y las metas mensuales de cinco sucursales. El director
quiere un reporte de texto, una línea por sucursal, con la venta, el
porcentaje de cumplimiento y una marca de alerta si quedó por debajo del
95\%.
\end{ejemplo}

\begin{rcode}
sucursales <- c("Centro", "Norte", "Sur", "Plaza A", "Outlet")
ventas_mes <- c(485000, 468900, 391200, 552300, 298700)
metas_mes <- c(520000, 450000, 420000, 540000, 280000)

cat("REPORTE DE CUMPLIMIENTO - MARZO 2023\n")
for (i in seq_along(sucursales)) {
  cumpl <- ventas_mes[i] / metas_mes[i] * 100
  alerta <- if (cumpl < 95) "  <- REVISAR" else ""
  cat(sprintf("%-8s $%9s  %6.1f%%%s\n", sucursales[i],
              format(ventas_mes[i], big.mark = ","), cumpl, alerta))
}
\end{rcode}
\begin{rsalida}
REPORTE DE CUMPLIMIENTO - MARZO 2023
Centro   $  485,000    93.3%  <- REVISAR
Norte    $  468,900   104.2%
Sur      $  391,200    93.1%  <- REVISAR
Plaza A  $  552,300   102.3%
Outlet   $  298,700   106.7%
\end{rsalida}

Vamos por partes:

\begin{itemize}
  \item \codigo{seq\_along(sucursales)} produce \codigo{1 2 3 4 5}, una
        posición por sucursal. Es más seguro que escribir \codigo{1:5}
        porque se ajusta solo si agregas sucursales.
  \item En cada vuelta, \codigo{i} vale 1, luego 2, etc., y con
        \codigo{[i]} tomas el dato de esa sucursal.
  \item \codigo{cat()} imprime texto tal cual, sin el \codigo{[1]} ni las
        comillas de \codigo{print()}. El \codigo{\textbackslash n} es un
        salto de línea.
  \item En \codigo{sprintf()}, \codigo{\%-8s} reserva 8 espacios para el
        texto alineado a la izquierda y \codigo{\%9s} reserva 9 alineados a
        la derecha: así las columnas quedan parejas.
\end{itemize}

\textbf{¿Qué nos dice esto?} Centro y Sur quedaron por debajo del 95\% de
su meta y hay que revisarlas; Outlet es la que más superó su meta
(106.7\%), aunque es la que menos vende en pesos.

\subsection{Repetir mientras se cumpla algo: \texorpdfstring{\codigo{while}}{while}}

\codigo{while} repite un bloque \emph{mientras} una condición sea
verdadera. Se usa cuando no sabes de antemano cuántas vueltas darás.

\begin{ejemplo}{¿Cuántos meses para llegar a la meta?}
Una tienda vende hoy \$350,000 al mes y su plan de negocio prevé crecer 3\%
mensual. ¿En cuántos meses alcanzará ventas mensuales de \$450,000?
\end{ejemplo}

\begin{rcode}
venta_actual <- 350000
crecimiento_mensual <- 0.03
meta_mensual <- 450000
meses <- 0

while (venta_actual < meta_mensual) {
  venta_actual <- venta_actual * (1 + crecimiento_mensual)
  meses <- meses + 1
}
meses
round(venta_actual, 2)
\end{rcode}
\begin{rsalida}
[1] 9
[1] 456670.6
\end{rsalida}

En cada vuelta la venta crece 3\% y el contador de meses sube en uno. El
bucle se detiene en cuanto la venta alcanza la meta: después de 9 meses la
tienda vendería unos \$456,671.

\begin{advertencia}
Si la condición de un \codigo{while} nunca se vuelve \codigo{FALSE} (por
ejemplo, con un crecimiento de 0\%), el bucle no termina jamás: es un
\termino{bucle infinito}. Si te pasa, presiona \tecla{Esc} o el botón rojo
\menu{Stop} de la consola.
\end{advertencia}

\subsection{En R se prefiere lo vectorizado}

Si vienes de Excel o de otro lenguaje, tu primer instinto será escribir
bucles para todo. En R casi siempre existe una forma vectorizada más corta,
más rápida y menos propensa a errores. Compara el bucle anterior de
cumplimiento con esta sola línea:

\begin{rcode}
round(ventas_mes / metas_mes * 100, 1)       # Cumplimiento de las 5
sucursales[ventas_mes / metas_mes < 0.95]    # Las que hay que revisar
\end{rcode}
\begin{rsalida}
[1]  93.3 104.2  93.1 102.3 106.7
[1] "Centro" "Sur"   
\end{rsalida}

Usa \codigo{for} cuando cada vuelta haga algo que no es un simple cálculo:
imprimir un reporte, leer un archivo distinto, guardar un Excel por
sucursal. Para aplicar una función a cada elemento de una lista existe
también la familia \codigo{apply}: \codigo{sapply()} aplica una función y
simplifica el resultado a un vector, y \codigo{vapply()} hace lo mismo
pero te obliga a declarar qué tipo de resultado esperas, lo que la hace más
segura.

\begin{rcode}
ventas_por_tienda <- list(
  Centro = c(15200, 13900, 16800),
  Norte = c(12100, 12900),
  Outlet = c(8400, 9100, 7600, 9900)
)
sapply(ventas_por_tienda, sum)                  # Total de cada tienda
sapply(ventas_por_tienda, length)               # Días con datos
vapply(ventas_por_tienda, mean, numeric(1))     # Promedio (seguro)
\end{rcode}
\begin{rsalida}
Centro  Norte Outlet 
 45900  25000  35000 
Centro  Norte Outlet 
     3      2      4 
Centro  Norte Outlet 
 15300  12500   8750 
\end{rsalida}

\codigo{numeric(1)} le dice a \codigo{vapply()}: ``cada resultado debe ser
un número''. Si alguna tienda devolviera texto o varios valores, R te
avisaría con un error en lugar de darte un resultado raro. En el
Módulo~\ref{cap:manipulacion} verás que con \paquete{dplyr} la mayoría de
estos cálculos por grupo se hacen sin bucles.

\begin{hazlotu}
Con los vectores \codigo{sucursales} y \codigo{ventas\_mes} usa
\codigo{ifelse()} para clasificar cada sucursal como ``Grande'' (ventas
$\geq$ \$450,000) o ``Chica''. Después escribe un \codigo{for} que imprima
con \codigo{cat()} una línea como \codigo{"Norte: Grande"} por sucursal.
\end{hazlotu}

% ============================================================================
\section{Tus propias funciones: una caja de herramientas BI}
\label{sec:m1-funciones}
% ============================================================================

Si copias y pegas el mismo cálculo tres veces, es momento de convertirlo en
una \termino{función}. Una función propia es como una fórmula personalizada
de Excel: le pones nombre una vez y la usas siempre que la necesites. Si
mañana cambia la regla de negocio, la corriges en un solo lugar.

\subsection{Anatomía de una función}

\begin{rcode}
# Calcula el margen de utilidad de uno o varios productos.
# Argumentos:
#   precio: precio de venta (número o vector)
#   costo:  costo unitario (número o vector)
#   en_pct: TRUE devuelve porcentaje (0-100); FALSE, proporción (0-1)
# Devuelve: el margen, redondeado a 1 decimal si es porcentaje
calcular_margen <- function(precio, costo, en_pct = TRUE) {
  margen <- (precio - costo) / precio
  if (en_pct) {
    return(round(margen * 100, 1))
  }
  margen
}
\end{rcode}

\begin{itemize}
  \item \codigo{calcular\_margen} es el nombre; \codigo{function()} crea la
        función.
  \item Entre paréntesis van los \termino{argumentos}: los datos que la
        función necesita. \codigo{en\_pct = TRUE} tiene un
        \termino{valor por defecto}: si no lo indicas, vale \codigo{TRUE}.
  \item Entre llaves va el \termino{cuerpo}: los cálculos.
  \item \codigo{return()} devuelve un resultado y termina la función. Si no
        lo usas, la función devuelve el resultado de la última línea
        (aquí, \codigo{margen}).
  \item Los comentarios de arriba \emph{documentan} la función: qué hace,
        qué recibe y qué devuelve. Es un hábito profesional que te salvará
        cuando vuelvas a este código en unos meses.
\end{itemize}

Definir la función no calcula nada: solo la ``guarda'' en el entorno.
Ahora úsala:

\begin{rcode}
calcular_margen(450, 280)                        # Posicional
calcular_margen(precio = 450, costo = 280, en_pct = FALSE)
calcular_margen(productos$precio, productos$costo)  # Vectorizada
\end{rcode}
\begin{rsalida}
[1] 37.8
[1] 0.3777778
[1] 25.0 65.6 25.4 41.0 31.4 57.3 46.9 44.7
\end{rsalida}

Como está construida con operaciones vectorizadas, la misma función sirve
para un producto o para todo el catálogo. Escribir los nombres de los
argumentos (\codigo{precio = 450}) es opcional, pero hace el código más
claro, sobre todo cuando hay argumentos con valor por defecto.

\subsection{Validar los datos: \texorpdfstring{\codigo{stop()}}{stop()} y \texorpdfstring{\codigo{warning()}}{warning()}}

Una buena función se protege de datos absurdos. \codigo{stop()} detiene la
función con un mensaje de error claro; \codigo{warning()} deja que continúe
pero avisa que algo es sospechoso.

\begin{rcode}
# ROI (retorno sobre la inversión) en porcentaje
calcular_roi <- function(ganancia, inversion) {
  if (!is.numeric(ganancia) || !is.numeric(inversion)) {
    stop("ganancia e inversion deben ser números")
  }
  if (any(inversion <= 0)) {
    stop("La inversión debe ser mayor que cero")
  }
  (ganancia - inversion) / inversion * 100
}

# Variación porcentual entre un periodo y el anterior
variacion_porcentual <- function(actual, anterior) {
  if (any(anterior == 0)) {
    warning("Hay periodos anteriores en 0; su variación queda como NA")
  }
  resultado <- (actual - anterior) / anterior * 100
  resultado[anterior == 0] <- NA
  round(resultado, 1)
}
\end{rcode}

El operador \codigo{||} es un ``o'' para un solo valor (dentro de
\codigo{if} se prefiere a \codigo{|}) y \codigo{any()} responde
\codigo{TRUE} si \emph{alguno} de los elementos cumple la condición. Veamos
qué pasa con datos inválidos:

\begin{rnoejecutar}
calcular_roi(5000, 0)
\end{rnoejecutar}
\begin{rsalida}
Error in calcular_roi(5000, 0) : La inversión debe ser mayor que cero
\end{rsalida}

\begin{rnoejecutar}
variacion_porcentual(c(120, 50), c(100, 0))
\end{rnoejecutar}
\begin{rsalida}
[1] 20 NA
Warning message:
In variacion_porcentual(c(120, 50), c(100, 0)) :
  Hay periodos anteriores en 0; su variación queda como NA
\end{rsalida}

Sin la validación, dividir entre cero daría \codigo{Inf} (infinito) y ese
valor absurdo se colaría en tu reporte. Con ella, el error explica
exactamente qué está mal.

\subsection{La caja de herramientas completa}

Agreguemos una función para segmentar clientes según su gasto anual y
usemos las cuatro juntas.

\begin{rcode}
# Segmenta clientes por gasto anual:
# Bronce < $5,000 <= Plata < $20,000 <= Oro < $50,000 <= Platino
clasificar_cliente <- function(gasto_anual,
                               cortes = c(5000, 20000, 50000)) {
  if (any(gasto_anual < 0, na.rm = TRUE)) {
    stop("El gasto no puede ser negativo")
  }
  ifelse(gasto_anual >= cortes[3], "Platino",
         ifelse(gasto_anual >= cortes[2], "Oro",
                ifelse(gasto_anual >= cortes[1], "Plata", "Bronce")))
}
\end{rcode}

\begin{ejemplo}{Usar la caja de herramientas}
Una campaña de marketing costó \$80,000 y generó \$116,000 de utilidad.
Las ventas trimestrales fueron \$126,274, \$120,717, \$134,427 y \$132,498.
Y seis clientes gastaron en el año \$3,200, \$18,500, \$52,000, \$7,800,
\$25,000 y \$950. Calcula el ROI, la variación trimestral y el segmento de
cada cliente.
\end{ejemplo}

\begin{rcode}
calcular_roi(ganancia = 116000, inversion = 80000)

ventas_trim <- c(Q1 = 126274, Q2 = 120717, Q3 = 134427, Q4 = 132498)
variacion_porcentual(ventas_trim[-1], ventas_trim[-4])

gastos <- c(3200, 18500, 52000, 7800, 25000, 950)
segmentos <- clasificar_cliente(gastos)
segmentos
table(factor(segmentos, levels = c("Bronce", "Plata", "Oro",
                                    "Platino")))
\end{rcode}
\begin{rsalida}
[1] 45
  Q2   Q3   Q4 
-4.4 11.4 -1.4 
[1] "Bronce"  "Plata"   "Platino" "Plata"   "Oro"     "Bronce" 

 Bronce   Plata     Oro Platino 
      2       2       1       1 
\end{rsalida}

\textbf{¿Qué nos dice esto?} La campaña tuvo un ROI de 45\%: por cada peso
invertido se recuperó el peso y 45 centavos más. Las ventas cayeron 4.4\% del
primer al segundo trimestre, repuntaron 11.4\% en el tercero y bajaron
ligeramente en el cuarto. Y la cartera tiene dos clientes Bronce, dos Plata
y un cliente en cada uno de los niveles Oro y Platino.

Observa el truco de \codigo{ventas\_trim[-1]} (todos menos el primero, es
decir Q2 a Q4) contra \codigo{ventas\_trim[-4]} (todos menos el último, Q1
a Q3): así comparas cada trimestre con el anterior sin bucles. Los nombres
del resultado son los del periodo actual.

\begin{consejo}
Guarda tus funciones de uso frecuente en un script aparte, por ejemplo
\archivo{funciones\_bi.R}, y cárgalas al inicio de cada análisis con
\codigo{source("funciones\_bi.R")}. Así construyes tu propia caja de
herramientas y todos tus reportes calculan los KPIs de la misma forma.
\end{consejo}

\begin{hazlotu}
Escribe una función \codigo{precio\_final(precio, descuento = 0, iva =
0.16)} que aplique primero el descuento (como proporción, 0.15 = 15\%) y
luego el IVA, y que se detenga con \codigo{stop()} si el descuento no está
entre 0 y 1. Pruébala con \codigo{precio\_final(18999, 0.15)}; debería dar
18,733.01.
\end{hazlotu}

% ============================================================================
\section{Paquetes: superpoderes para R}
\label{sec:m1-paquetes}
% ============================================================================

R ``de fábrica'' (llamado \termino{R base}) ya hace muchísimo, pero su
verdadero poder está en los \termino{paquetes}: colecciones de funciones,
datos y documentación que otras personas escribieron y comparten gratis. Si
R fuera un teléfono, los paquetes serían las apps. Hay paquetes para leer
Excel (\paquete{readxl}), manipular datos (\paquete{dplyr}), hacer gráficas
(\paquete{ggplot2}), crear tableros web (\paquete{shiny}) o entrenar
modelos de \emph{machine learning}.

\subsection{Instalar y cargar}

La mayoría de los paquetes viven en \termino{CRAN} (\emph{Comprehensive R
Archive Network}), el repositorio oficial de R, con más de 20,000
paquetes revisados. Trabajar con un paquete tiene dos pasos, que conviene
no confundir:

\begin{enumerate}
  \item \textbf{Instalar} (una sola vez por computadora): descarga el
        paquete de internet. Es como descargar una app.
  \item \textbf{Cargar} (en cada sesión en la que lo uses): lo activa para
        la sesión actual. Es como abrir la app.
\end{enumerate}

\begin{rnoejecutar}
# Instalar (solo una vez; necesita internet). Nota las comillas.
install.packages("readxl")
install.packages(c("tidyverse", "writexl", "scales"))  # Varios a la vez
\end{rnoejecutar}

También puedes instalar desde el panel \menu{Packages > Install}. No pongas
\codigo{install.packages()} dentro de tus scripts de análisis: se volvería a
descargar el paquete cada vez que los ejecutes. Instala desde la consola y
deja en el script solo las líneas \codigo{library()}.

\begin{rcode}
library(readxl)                        # Cargar (en cada sesión)
library(writexl)
"readxl" %in% rownames(installed.packages())  # ¿Está instalado?
as.character(packageVersion("readxl"))  # ¿Qué versión tengo?
\end{rcode}
\begin{rsalida}
[1] TRUE
[1] "1.4.3"
\end{rsalida}

\begin{concepto}{Instalar contra cargar}
\codigo{install.packages("readxl")} va con comillas y se hace una vez.
\codigo{library(readxl)} va sin comillas y se escribe al inicio de cada
script que use el paquete. Si ves el error \codigo{could not find function
"read\_excel"}, casi siempre olvidaste el \codigo{library()}; si ves
\codigo{there is no package called 'readxl'}, te falta instalarlo.
\end{concepto}

\subsection{La notación \texorpdfstring{\codigo{paquete::funcion()}}{paquete::funcion()}}

Ya usaste \codigo{scales::dollar()} y \codigo{tibble::tibble()}. Los dos
signos de dos puntos permiten usar una función de un paquete
\emph{instalado} sin cargarlo completo con \codigo{library()}. Es útil
cuando solo necesitas una función, y además deja claro de dónde viene:

\begin{rcode}
scales::dollar(513916.6)      # Sin library(scales)
\end{rcode}
\begin{rsalida}
[1] "$513,917"
\end{rsalida}

\subsection{El tidyverse y su mensaje de bienvenida}

El \termino{tidyverse} es una colección de paquetes diseñados para
trabajar juntos en ciencia de datos: \paquete{readr} (leer archivos),
\paquete{dplyr} (manipular tablas), \paquete{tidyr} (reorganizarlas),
\paquete{ggplot2} (gráficas), \paquete{stringr} (texto),
\paquete{lubridate} (fechas), \paquete{forcats} (factores),
\paquete{tibble} y \paquete{purrr}. Con una sola línea los cargas todos:

\begin{rcode}
library(tidyverse)
\end{rcode}

Al ejecutarla, la consola muestra un mensaje de bienvenida en dos partes.
No es un error; conviene entenderlo:

\begin{itemize}
  \item \textbf{Attaching core tidyverse packages}: la lista de los nueve
        paquetes que se cargaron, cada uno con su versión y una palomita.
  \item \textbf{Conflicts}: avisa de funciones con el mismo nombre en dos
        paquetes. Verás líneas como \codigo{dplyr::filter() masks
        stats::filter()}, que significan: ``el paquete \paquete{stats} (que
        viene con R) ya tenía una función \codigo{filter()}; a partir de
        ahora, cuando escribas \codigo{filter()} se usará la de
        \paquete{dplyr}''. Lo mismo ocurre con \codigo{lag()}.
\end{itemize}

Estos conflictos son normales y esperados: la versión de \paquete{dplyr} es
justo la que quieres. Si alguna vez necesitas la otra, escribe el nombre
completo, por ejemplo \codigo{stats::filter()}. El mensaje aparece solo al
cargar el paquete; en este libro no lo reproducimos para ahorrar espacio.

\begin{hazlotu}
En la pestaña \menu{Packages} de RStudio busca \paquete{readxl} y haz clic
en su nombre: se abrirá su documentación en \menu{Help}. Después ejecuta
\codigo{packageVersion("dplyr")} y
\codigo{"janitor" \%in\% rownames(installed.packages())} para saber si
tienes instalado el paquete \paquete{janitor}.
\end{hazlotu}

% --- CONTINUARÁ ---
